% Ngân hàng đề Toán 12 (KNTT) -- Chương II, Bài 8, Dạng 08
% Dạng: Tích vô hướng, góc, vuông góc bằng tọa độ
%% Phương án ĐẦU TIÊN trong \cauTN luôn là đáp án đúng (gói trondethi tự đảo).

\cauTN[id=C2-B8-D08-NB-001,muc=NB,loigiai={Ta có $\vec{a}\cdot\vec{b}=3\cdot3+1\cdot-2+2\cdot-3=1$.}]{Trong không gian $Oxyz$, cho $\vec{a}=(3;1;2)$, $\vec{b}=(3;-2;-3)$. Tích vô hướng $\vec{a}\cdot\vec{b}$ bằng}
  {$1$.}
  {$3$.}{$2$.}{$-2$.}

\cauTN[id=C2-B8-D08-NB-002,muc=NB,loigiai={Ta có $\vec{a}\cdot\vec{b}=2\cdot1-3\cdot-1-2\cdot0=5$.}]{Trong không gian $Oxyz$, cho $\vec{a}=(2;-3;-2)$, $\vec{b}=(1;-1;0)$. Tích vô hướng $\vec{a}\cdot\vec{b}$ bằng}
  {$5$.}
  {$-4$.}{$6$.}{$3$.}

\cauTN[id=C2-B8-D08-NB-003,muc=NB,loigiai={Ta có $\vec{a}\cdot\vec{b}=-1\cdot2+2\cdot1+3\cdot-1=-3$.}]{Trong không gian $Oxyz$, cho $\vec{a}=(-1;2;3)$, $\vec{b}=(2;1;-1)$. Tích vô hướng $\vec{a}\cdot\vec{b}$ bằng}
  {$-3$.}
  {$-6$.}{$-5$.}{$0$.}

\cauTN[id=C2-B8-D08-NB-004,muc=NB,loigiai={Ta có $\vec{a}\cdot\vec{b}=-3\cdot2+1\cdot-3-3\cdot1=-12$.}]{Trong không gian $Oxyz$, cho $\vec{a}=(-3;1;-3)$, $\vec{b}=(2;-3;1)$. Tích vô hướng $\vec{a}\cdot\vec{b}$ bằng}
  {$-12$.}
  {$-10$.}{$-13$.}{$-6$.}

\cauTN[id=C2-B8-D08-NB-005,muc=NB,loigiai={Ta có $\vec{a}\cdot\vec{b}=-1\cdot3+0\cdot2+0\cdot1=-3$.}]{Trong không gian $Oxyz$, cho $\vec{a}=(-1;0;0)$, $\vec{b}=(3;2;1)$. Tích vô hướng $\vec{a}\cdot\vec{b}$ bằng}
  {$-3$.}
  {$3$.}{$-9$.}{$-6$.}

\cauTN[id=C2-B8-D08-NB-006,muc=NB,loigiai={Ta có $\vec{a}\cdot\vec{b}=2\cdot1+1\cdot-2-1\cdot1=-1$.}]{Trong không gian $Oxyz$, cho $\vec{a}=(2;1;-1)$, $\vec{b}=(1;-2;1)$. Tích vô hướng $\vec{a}\cdot\vec{b}$ bằng}
  {$-1$.}
  {$-3$.}{$0$.}{$-2$.}

\cauTN[id=C2-B8-D08-NB-007,muc=NB,loigiai={Ta có $\vec{a}\cdot\vec{b}=1\cdot-1+2\cdot-2-1\cdot0=-5$.}]{Trong không gian $Oxyz$, cho $\vec{a}=(1;2;-1)$, $\vec{b}=(-1;-2;0)$. Tích vô hướng $\vec{a}\cdot\vec{b}$ bằng}
  {$-5$.}
  {$-7$.}{$-3$.}{$-10$.}

\cauTN[id=C2-B8-D08-NB-008,muc=NB,loigiai={Ta có $\vec{a}\cdot\vec{b}=1\cdot1-3\cdot-2-2\cdot-2=11$.}]{Trong không gian $Oxyz$, cho $\vec{a}=(1;-3;-2)$, $\vec{b}=(1;-2;-2)$. Tích vô hướng $\vec{a}\cdot\vec{b}$ bằng}
  {$11$.}
  {$33$.}{$14$.}{$13$.}

\cauTN[id=C2-B8-D08-NB-009,muc=NB,loigiai={Ta có $\vec{a}\cdot\vec{b}=-3\cdot-1+0\cdot1-1\cdot0=3$.}]{Trong không gian $Oxyz$, cho $\vec{a}=(-3;0;-1)$, $\vec{b}=(-1;1;0)$. Tích vô hướng $\vec{a}\cdot\vec{b}$ bằng}
  {$3$.}
  {$6$.}{$2$.}{$-2$.}

\cauTN[id=C2-B8-D08-NB-010,muc=NB,loigiai={Ta có $\vec{a}\cdot\vec{b}=3\cdot1+2\cdot3-1\cdot-3=12$.}]{Trong không gian $Oxyz$, cho $\vec{a}=(3;2;-1)$, $\vec{b}=(1;3;-3)$. Tích vô hướng $\vec{a}\cdot\vec{b}$ bằng}
  {$12$.}
  {$2$.}{$24$.}{$6$.}

\cauTN[id=C2-B8-D08-NB-011,muc=NB,loigiai={Ta có $\vec{a}\cdot\vec{b}=3\cdot1+3\cdot-1-1\cdot2=-2$.}]{Trong không gian $Oxyz$, cho $\vec{a}=(3;3;-1)$, $\vec{b}=(1;-1;2)$. Tích vô hướng $\vec{a}\cdot\vec{b}$ bằng}
  {$-2$.}
  {$2$.}{$0$.}{$-4$.}

\cauTN[id=C2-B8-D08-NB-012,muc=NB,loigiai={Ta có $\vec{a}\cdot\vec{b}=0\cdot2+0\cdot-2-3\cdot-3=9$.}]{Trong không gian $Oxyz$, cho $\vec{a}=(0;0;-3)$, $\vec{b}=(2;-2;-3)$. Tích vô hướng $\vec{a}\cdot\vec{b}$ bằng}
  {$9$.}
  {$27$.}{$18$.}{$12$.}

\cauTN[id=C2-B8-D08-NB-013,muc=NB,loigiai={Ta có $\cos(\vec{a},\vec{b})=\dfrac{\vec{a}\cdot\vec{b}}{|\vec{a}|\,|\vec{b}|}=\dfrac{0}{\sqrt{17}\cdot\sqrt{2}}=0$, suy ra góc bằng $90^\circ$.}]{Trong không gian $Oxyz$, cho $\vec{a}=(-2;-2;-3)$, $\vec{b}=(-1;1;0)$. Góc giữa hai vectơ $\vec{a}$ và $\vec{b}$ bằng}
  {$90^\circ$.}
  {$60^\circ$.}{$45^\circ$.}{$120^\circ$.}

\cauTN[id=C2-B8-D08-NB-014,muc=NB,loigiai={Ta có $\cos(\vec{a},\vec{b})=\dfrac{\vec{a}\cdot\vec{b}}{|\vec{a}|\,|\vec{b}|}=\dfrac{0}{\sqrt{12}\cdot\sqrt{6}}=0$, suy ra góc bằng $90^\circ$.}]{Trong không gian $Oxyz$, cho $\vec{a}=(-2;-2;2)$, $\vec{b}=(-1;2;1)$. Góc giữa hai vectơ $\vec{a}$ và $\vec{b}$ bằng}
  {$90^\circ$.}
  {$135^\circ$.}{$150^\circ$.}{$45^\circ$.}

\cauTN[id=C2-B8-D08-NB-015,muc=NB,loigiai={Ta có $\cos(\vec{a},\vec{b})=\dfrac{\vec{a}\cdot\vec{b}}{|\vec{a}|\,|\vec{b}|}=\dfrac{0}{\sqrt{27}\cdot\sqrt{8}}=0$, suy ra góc bằng $90^\circ$.}]{Trong không gian $Oxyz$, cho $\vec{a}=(3;-3;-3)$, $\vec{b}=(-2;-2;0)$. Góc giữa hai vectơ $\vec{a}$ và $\vec{b}$ bằng}
  {$90^\circ$.}
  {$120^\circ$.}{$30^\circ$.}{$135^\circ$.}

\cauTN[id=C2-B8-D08-NB-016,muc=NB,loigiai={Ta có $\cos(\vec{a},\vec{b})=\dfrac{\vec{a}\cdot\vec{b}}{|\vec{a}|\,|\vec{b}|}=\dfrac{0}{\sqrt{18}\cdot\sqrt{27}}=0$, suy ra góc bằng $90^\circ$.}]{Trong không gian $Oxyz$, cho $\vec{a}=(3;3;0)$, $\vec{b}=(3;-3;-3)$. Góc giữa hai vectơ $\vec{a}$ và $\vec{b}$ bằng}
  {$90^\circ$.}
  {$135^\circ$.}{$45^\circ$.}{$30^\circ$.}

\cauTN[id=C2-B8-D08-NB-017,muc=NB,loigiai={Ta có $\cos(\vec{a},\vec{b})=\dfrac{\vec{a}\cdot\vec{b}}{|\vec{a}|\,|\vec{b}|}=\dfrac{0}{\sqrt{27}\cdot\sqrt{2}}=0$, suy ra góc bằng $90^\circ$.}]{Trong không gian $Oxyz$, cho $\vec{a}=(-3;3;3)$, $\vec{b}=(0;1;-1)$. Góc giữa hai vectơ $\vec{a}$ và $\vec{b}$ bằng}
  {$90^\circ$.}
  {$120^\circ$.}{$60^\circ$.}{$30^\circ$.}

\cauTN[id=C2-B8-D08-NB-018,muc=NB,loigiai={Ta có $\cos(\vec{a},\vec{b})=\dfrac{\vec{a}\cdot\vec{b}}{|\vec{a}|\,|\vec{b}|}=\dfrac{0}{\sqrt{10}\cdot\sqrt{14}}=0$, suy ra góc bằng $90^\circ$.}]{Trong không gian $Oxyz$, cho $\vec{a}=(3;0;1)$, $\vec{b}=(1;-2;-3)$. Góc giữa hai vectơ $\vec{a}$ và $\vec{b}$ bằng}
  {$90^\circ$.}
  {$60^\circ$.}{$45^\circ$.}{$150^\circ$.}

\cauTN[id=C2-B8-D08-NB-019,muc=NB,loigiai={Ta có $\cos(\vec{a},\vec{b})=\dfrac{\vec{a}\cdot\vec{b}}{|\vec{a}|\,|\vec{b}|}=\dfrac{-1}{\sqrt{2}\cdot\sqrt{2}}=- \dfrac{1}{2}$, suy ra góc bằng $120^\circ$.}]{Trong không gian $Oxyz$, cho $\vec{a}=(1;0;-1)$, $\vec{b}=(-1;-1;0)$. Góc giữa hai vectơ $\vec{a}$ và $\vec{b}$ bằng}
  {$120^\circ$.}
  {$135^\circ$.}{$45^\circ$.}{$60^\circ$.}

\cauTN[id=C2-B8-D08-NB-020,muc=NB,loigiai={Ta có $\cos(\vec{a},\vec{b})=\dfrac{\vec{a}\cdot\vec{b}}{|\vec{a}|\,|\vec{b}|}=\dfrac{0}{\sqrt{17}\cdot\sqrt{13}}=0$, suy ra góc bằng $90^\circ$.}]{Trong không gian $Oxyz$, cho $\vec{a}=(2;-2;3)$, $\vec{b}=(0;3;2)$. Góc giữa hai vectơ $\vec{a}$ và $\vec{b}$ bằng}
  {$90^\circ$.}
  {$60^\circ$.}{$30^\circ$.}{$120^\circ$.}

\cauTN[id=C2-B8-D08-NB-021,muc=NB,loigiai={Ta có $\cos(\vec{a},\vec{b})=\dfrac{\vec{a}\cdot\vec{b}}{|\vec{a}|\,|\vec{b}|}=\dfrac{0}{\sqrt{14}\cdot\sqrt{10}}=0$, suy ra góc bằng $90^\circ$.}]{Trong không gian $Oxyz$, cho $\vec{a}=(-1;-2;3)$, $\vec{b}=(-3;0;-1)$. Góc giữa hai vectơ $\vec{a}$ và $\vec{b}$ bằng}
  {$90^\circ$.}
  {$45^\circ$.}{$30^\circ$.}{$150^\circ$.}

\cauTN[id=C2-B8-D08-NB-022,muc=NB,loigiai={Ta có $\cos(\vec{a},\vec{b})=\dfrac{\vec{a}\cdot\vec{b}}{|\vec{a}|\,|\vec{b}|}=\dfrac{0}{\sqrt{6}\cdot\sqrt{18}}=0$, suy ra góc bằng $90^\circ$.}]{Trong không gian $Oxyz$, cho $\vec{a}=(1;-2;-1)$, $\vec{b}=(3;0;3)$. Góc giữa hai vectơ $\vec{a}$ và $\vec{b}$ bằng}
  {$90^\circ$.}
  {$150^\circ$.}{$120^\circ$.}{$135^\circ$.}

\cauTN[id=C2-B8-D08-NB-023,muc=NB,loigiai={Hai vectơ vuông góc khi tích vô hướng bằng $0$: $\vec{a}\cdot\vec{d}=0=0$. Các cặp còn lại có tích vô hướng khác $0$.}]{Trong không gian $Oxyz$, cho $\vec{a}=(-2;2;0)$, $\vec{b}=(3;-1;0)$, $\vec{c}=(0;3;-2)$, $\vec{d}=(3;3;1)$. Cặp vectơ nào dưới đây vuông góc với nhau?}
  {$\vec{a}\text{ và }\vec{d}$.}
  {$\vec{b}\text{ và }\vec{d}$.}{$\vec{a}\text{ và }\vec{c}$.}{$\vec{a}\text{ và }\vec{b}$.}

\cauTN[id=C2-B8-D08-NB-024,muc=NB,loigiai={Hai vectơ vuông góc khi tích vô hướng bằng $0$: $\vec{c}\cdot\vec{d}=0=0$. Các cặp còn lại có tích vô hướng khác $0$.}]{Trong không gian $Oxyz$, cho $\vec{a}=(-2;-1;-1)$, $\vec{b}=(-2;1;0)$, $\vec{c}=(1;3;2)$, $\vec{d}=(-1;1;-1)$. Cặp vectơ nào dưới đây vuông góc với nhau?}
  {$\vec{c}\text{ và }\vec{d}$.}
  {$\vec{b}\text{ và }\vec{d}$.}{$\vec{a}\text{ và }\vec{b}$.}{$\vec{a}\text{ và }\vec{c}$.}

\cauTN[id=C2-B8-D08-NB-025,muc=NB,loigiai={Hai vectơ vuông góc khi tích vô hướng bằng $0$: $\vec{c}\cdot\vec{d}=0=0$. Các cặp còn lại có tích vô hướng khác $0$.}]{Trong không gian $Oxyz$, cho $\vec{a}=(1;1;2)$, $\vec{b}=(1;-2;3)$, $\vec{c}=(0;-2;0)$, $\vec{d}=(0;0;-1)$. Cặp vectơ nào dưới đây vuông góc với nhau?}
  {$\vec{c}\text{ và }\vec{d}$.}
  {$\vec{b}\text{ và }\vec{d}$.}{$\vec{a}\text{ và }\vec{c}$.}{$\vec{a}\text{ và }\vec{b}$.}

\cauTN[id=C2-B8-D08-NB-026,muc=NB,loigiai={Hai vectơ vuông góc khi tích vô hướng bằng $0$: $\vec{a}\cdot\vec{b}=0=0$. Các cặp còn lại có tích vô hướng khác $0$.}]{Trong không gian $Oxyz$, cho $\vec{a}=(-2;2;-1)$, $\vec{b}=(-3;-3;0)$, $\vec{c}=(3;-1;1)$, $\vec{d}=(0;-2;2)$. Cặp vectơ nào dưới đây vuông góc với nhau?}
  {$\vec{a}\text{ và }\vec{b}$.}
  {$\vec{a}\text{ và }\vec{d}$.}{$\vec{b}\text{ và }\vec{d}$.}{$\vec{a}\text{ và }\vec{c}$.}

\cauTN[id=C2-B8-D08-NB-027,muc=NB,loigiai={Hai vectơ vuông góc khi tích vô hướng bằng $0$: $\vec{b}\cdot\vec{c}=0=0$. Các cặp còn lại có tích vô hướng khác $0$.}]{Trong không gian $Oxyz$, cho $\vec{a}=(-2;-1;-2)$, $\vec{b}=(-1;-1;0)$, $\vec{c}=(0;0;-2)$, $\vec{d}=(3;0;2)$. Cặp vectơ nào dưới đây vuông góc với nhau?}
  {$\vec{b}\text{ và }\vec{c}$.}
  {$\vec{c}\text{ và }\vec{d}$.}{$\vec{a}\text{ và }\vec{b}$.}{$\vec{a}\text{ và }\vec{c}$.}

\cauTN[id=C2-B8-D08-NB-028,muc=NB,loigiai={Hai vectơ vuông góc khi tích vô hướng bằng $0$: $\vec{a}\cdot\vec{d}=0=0$. Các cặp còn lại có tích vô hướng khác $0$.}]{Trong không gian $Oxyz$, cho $\vec{a}=(-1;1;-1)$, $\vec{b}=(1;2;-1)$, $\vec{c}=(2;3;-3)$, $\vec{d}=(-2;-3;-1)$. Cặp vectơ nào dưới đây vuông góc với nhau?}
  {$\vec{a}\text{ và }\vec{d}$.}
  {$\vec{b}\text{ và }\vec{d}$.}{$\vec{c}\text{ và }\vec{d}$.}{$\vec{a}\text{ và }\vec{b}$.}

\cauTN[id=C2-B8-D08-NB-029,muc=NB,loigiai={Hai vectơ vuông góc khi tích vô hướng bằng $0$: $\vec{c}\cdot\vec{d}=0=0$. Các cặp còn lại có tích vô hướng khác $0$.}]{Trong không gian $Oxyz$, cho $\vec{a}=(-2;3;-3)$, $\vec{b}=(1;0;-1)$, $\vec{c}=(0;3;-1)$, $\vec{d}=(3;0;0)$. Cặp vectơ nào dưới đây vuông góc với nhau?}
  {$\vec{c}\text{ và }\vec{d}$.}
  {$\vec{a}\text{ và }\vec{d}$.}{$\vec{b}\text{ và }\vec{c}$.}{$\vec{a}\text{ và }\vec{c}$.}

\cauTN[id=C2-B8-D08-NB-030,muc=NB,loigiai={Hai vectơ vuông góc khi tích vô hướng bằng $0$: $\vec{a}\cdot\vec{d}=0=0$. Các cặp còn lại có tích vô hướng khác $0$.}]{Trong không gian $Oxyz$, cho $\vec{a}=(3;-2;-2)$, $\vec{b}=(-3;0;2)$, $\vec{c}=(-2;3;0)$, $\vec{d}=(-2;-3;0)$. Cặp vectơ nào dưới đây vuông góc với nhau?}
  {$\vec{a}\text{ và }\vec{d}$.}
  {$\vec{a}\text{ và }\vec{c}$.}{$\vec{b}\text{ và }\vec{d}$.}{$\vec{a}\text{ và }\vec{b}$.}

\cauTN[id=C2-B8-D08-NB-031,muc=NB,loigiai={Hai vectơ vuông góc khi tích vô hướng bằng $0$: $\vec{a}\cdot\vec{b}=0=0$. Các cặp còn lại có tích vô hướng khác $0$.}]{Trong không gian $Oxyz$, cho $\vec{a}=(3;2;2)$, $\vec{b}=(2;-1;-2)$, $\vec{c}=(1;-1;3)$, $\vec{d}=(0;3;0)$. Cặp vectơ nào dưới đây vuông góc với nhau?}
  {$\vec{a}\text{ và }\vec{b}$.}
  {$\vec{a}\text{ và }\vec{c}$.}{$\vec{b}\text{ và }\vec{c}$.}{$\vec{c}\text{ và }\vec{d}$.}

\cauTN[id=C2-B8-D08-NB-032,muc=NB,loigiai={Hai vectơ vuông góc khi tích vô hướng bằng $0$: $\vec{c}\cdot\vec{d}=0=0$. Các cặp còn lại có tích vô hướng khác $0$.}]{Trong không gian $Oxyz$, cho $\vec{a}=(0;-1;-3)$, $\vec{b}=(-3;2;1)$, $\vec{c}=(-3;-1;-1)$, $\vec{d}=(0;2;-2)$. Cặp vectơ nào dưới đây vuông góc với nhau?}
  {$\vec{c}\text{ và }\vec{d}$.}
  {$\vec{b}\text{ và }\vec{d}$.}{$\vec{a}\text{ và }\vec{c}$.}{$\vec{a}\text{ và }\vec{b}$.}

\cauTN[id=C2-B8-D08-NB-033,muc=NB,loigiai={Hai vectơ vuông góc khi tích vô hướng bằng $0$: $\vec{c}\cdot\vec{d}=0=0$. Các cặp còn lại có tích vô hướng khác $0$.}]{Trong không gian $Oxyz$, cho $\vec{a}=(-3;1;-2)$, $\vec{b}=(-2;-3;1)$, $\vec{c}=(3;-2;2)$, $\vec{d}=(2;1;-2)$. Cặp vectơ nào dưới đây vuông góc với nhau?}
  {$\vec{c}\text{ và }\vec{d}$.}
  {$\vec{b}\text{ và }\vec{c}$.}{$\vec{a}\text{ và }\vec{c}$.}{$\vec{b}\text{ và }\vec{d}$.}

\cauTN[id=C2-B8-D08-NB-034,muc=NB,loigiai={Hai vectơ vuông góc khi tích vô hướng bằng $0$: $\vec{c}\cdot\vec{d}=0=0$. Các cặp còn lại có tích vô hướng khác $0$.}]{Trong không gian $Oxyz$, cho $\vec{a}=(3;-3;2)$, $\vec{b}=(0;2;0)$, $\vec{c}=(1;3;1)$, $\vec{d}=(-2;1;-1)$. Cặp vectơ nào dưới đây vuông góc với nhau?}
  {$\vec{c}\text{ và }\vec{d}$.}
  {$\vec{a}\text{ và }\vec{c}$.}{$\vec{b}\text{ và }\vec{d}$.}{$\vec{b}\text{ và }\vec{c}$.}

\cauTN[id=C2-B8-D08-NB-035,muc=NB,loigiai={Hai vectơ vuông góc khi tích vô hướng bằng $0$: $\vec{a}\cdot\vec{c}=0=0$. Các cặp còn lại có tích vô hướng khác $0$.}]{Trong không gian $Oxyz$, cho $\vec{a}=(-3;2;0)$, $\vec{b}=(3;2;3)$, $\vec{c}=(-2;-3;-3)$, $\vec{d}=(-3;-3;-1)$. Cặp vectơ nào dưới đây vuông góc với nhau?}
  {$\vec{a}\text{ và }\vec{c}$.}
  {$\vec{a}\text{ và }\vec{d}$.}{$\vec{a}\text{ và }\vec{b}$.}{$\vec{c}\text{ và }\vec{d}$.}

\cauTN[id=C2-B8-D08-NB-036,muc=NB,loigiai={Hai vectơ vuông góc khi tích vô hướng bằng $0$: $\vec{a}\cdot\vec{b}=0=0$. Các cặp còn lại có tích vô hướng khác $0$.}]{Trong không gian $Oxyz$, cho $\vec{a}=(3;1;0)$, $\vec{b}=(1;-3;-2)$, $\vec{c}=(-2;2;1)$, $\vec{d}=(0;2;3)$. Cặp vectơ nào dưới đây vuông góc với nhau?}
  {$\vec{a}\text{ và }\vec{b}$.}
  {$\vec{b}\text{ và }\vec{d}$.}{$\vec{b}\text{ và }\vec{c}$.}{$\vec{a}\text{ và }\vec{c}$.}

\cauDS[id=C2-B8-D08-NB-037,muc=NB,loigiai={Kiểm tra từng mệnh đề bằng công thức tọa độ. a) Đúng. b) Đúng. c) Sai. d) Sai.}]{Trong không gian $Oxyz$, cho $\vec{a}=(1;-1;-1)$, $\vec{b}=(-3;-2;-1)$. Xét tính đúng sai của các mệnh đề sau:}{%
  \y{$\vec{a}\cdot\vec{b}=0$.}{d}
  \y{Độ dài của vectơ $\vec{a}$ bằng $\sqrt{3}$.}{d}
  \y{$\vec{a}\cdot\vec{b}=-1$.}{s}
  \y{Góc giữa hai vectơ $\vec{a}$ và $\vec{b}$ bằng $45^\circ$.}{s}}

\cauDS[id=C2-B8-D08-NB-038,muc=NB,loigiai={Kiểm tra từng mệnh đề bằng công thức tọa độ. a) Sai. b) Sai. c) Đúng. d) Đúng.}]{Trong không gian $Oxyz$, cho $\vec{a}=(1;0;-2)$, $\vec{b}=(1;0;3)$. Xét tính đúng sai của các mệnh đề sau:}{%
  \y{$\vec{a}\cdot\vec{b}=-2$.}{s}
  \y{Góc giữa hai vectơ $\vec{a}$ và $\vec{b}$ bằng $90^\circ$.}{s}
  \y{Độ dài của vectơ $\vec{a}$ bằng $\sqrt{5}$.}{d}
  \y{Góc giữa hai vectơ $\vec{a}$ và $\vec{b}$ bằng $135^\circ$.}{d}}

\cauDS[id=C2-B8-D08-NB-039,muc=NB,loigiai={Kiểm tra từng mệnh đề bằng công thức tọa độ. a) Sai. b) Sai. c) Sai. d) Đúng.}]{Trong không gian $Oxyz$, cho $\vec{a}=(0;-1;1)$, $\vec{b}=(-3;-1;-1)$. Xét tính đúng sai của các mệnh đề sau:}{%
  \y{Độ dài của vectơ $\vec{a}$ bằng $\sqrt{3}$.}{s}
  \y{Độ dài của vectơ $\vec{a}$ bằng $1$.}{s}
  \y{$\vec{a}\cdot\vec{b}=-2$.}{s}
  \y{Góc giữa hai vectơ $\vec{a}$ và $\vec{b}$ bằng $90^\circ$.}{d}}

\cauDS[id=C2-B8-D08-NB-040,muc=NB,loigiai={Kiểm tra từng mệnh đề bằng công thức tọa độ. a) Sai. b) Đúng. c) Đúng. d) Đúng.}]{Trong không gian $Oxyz$, cho $\vec{a}=(1;-3;-1)$, $\vec{b}=(1;1;-2)$. Xét tính đúng sai của các mệnh đề sau:}{%
  \y{$\vec{a}\cdot\vec{b}=-4$.}{s}
  \y{Độ dài của vectơ $\vec{a}$ bằng $\sqrt{11}$.}{d}
  \y{Góc giữa hai vectơ $\vec{a}$ và $\vec{b}$ bằng $90^\circ$.}{d}
  \y{$\vec{a}\cdot\vec{b}=0$.}{d}}

\cauDS[id=C2-B8-D08-NB-041,muc=NB,loigiai={Kiểm tra từng mệnh đề bằng công thức tọa độ. a) Sai. b) Đúng. c) Sai. d) Đúng.}]{Trong không gian $Oxyz$, cho $\vec{a}=(2;1;0)$, $\vec{b}=(-1;-3;0)$. Xét tính đúng sai của các mệnh đề sau:}{%
  \y{$\vec{a}\cdot\vec{b}=-2$.}{s}
  \y{Góc giữa hai vectơ $\vec{a}$ và $\vec{b}$ bằng $135^\circ$.}{d}
  \y{Độ dài của vectơ $\vec{a}$ bằng $\sqrt{6}$.}{s}
  \y{$\vec{a}\cdot\vec{b}=-5$.}{d}}

\cauDS[id=C2-B8-D08-NB-042,muc=NB,loigiai={Kiểm tra từng mệnh đề bằng công thức tọa độ. a) Đúng. b) Sai. c) Sai. d) Sai.}]{Trong không gian $Oxyz$, cho $\vec{a}=(-1;-2;-3)$, $\vec{b}=(-3;3;-1)$. Xét tính đúng sai của các mệnh đề sau:}{%
  \y{Góc giữa hai vectơ $\vec{a}$ và $\vec{b}$ bằng $90^\circ$.}{d}
  \y{Độ dài của vectơ $\vec{a}$ bằng $\sqrt{15}$.}{s}
  \y{$\vec{a}\cdot\vec{b}=-2$.}{s}
  \y{Góc giữa hai vectơ $\vec{a}$ và $\vec{b}$ bằng $30^\circ$.}{s}}

\cauDS[id=C2-B8-D08-NB-043,muc=NB,loigiai={Kiểm tra từng mệnh đề bằng công thức tọa độ. a) Đúng. b) Sai. c) Sai. d) Đúng.}]{Trong không gian $Oxyz$, cho $\vec{a}=(0;1;-2)$, $\vec{b}=(0;-2;-1)$. Xét tính đúng sai của các mệnh đề sau:}{%
  \y{Độ dài của vectơ $\vec{a}$ bằng $\sqrt{5}$.}{d}
  \y{Góc giữa hai vectơ $\vec{a}$ và $\vec{b}$ bằng $60^\circ$.}{s}
  \y{Góc giữa hai vectơ $\vec{a}$ và $\vec{b}$ bằng $150^\circ$.}{s}
  \y{$\vec{a}\cdot\vec{b}=0$.}{d}}

\cauDS[id=C2-B8-D08-NB-044,muc=NB,loigiai={Kiểm tra từng mệnh đề bằng công thức tọa độ. a) Đúng. b) Đúng. c) Sai. d) Sai.}]{Trong không gian $Oxyz$, cho $\vec{a}=(-3;0;-3)$, $\vec{b}=(1;-1;-1)$. Xét tính đúng sai của các mệnh đề sau:}{%
  \y{Độ dài của vectơ $\vec{a}$ bằng $3\sqrt{2}$.}{d}
  \y{$\vec{a}\cdot\vec{b}=0$.}{d}
  \y{Độ dài của vectơ $\vec{a}$ bằng $\sqrt{22}$.}{s}
  \y{Góc giữa hai vectơ $\vec{a}$ và $\vec{b}$ bằng $150^\circ$.}{s}}

\cauDS[id=C2-B8-D08-NB-045,muc=NB,loigiai={Kiểm tra từng mệnh đề bằng công thức tọa độ. a) Đúng. b) Sai. c) Sai. d) Đúng.}]{Trong không gian $Oxyz$, cho $\vec{a}=(1;-3;2)$, $\vec{b}=(-3;-1;0)$. Xét tính đúng sai của các mệnh đề sau:}{%
  \y{$\vec{a}\cdot\vec{b}=0$.}{d}
  \y{Độ dài của vectơ $\vec{a}$ bằng $\sqrt{15}$.}{s}
  \y{Độ dài của vectơ $\vec{a}$ bằng $\sqrt{13}$.}{s}
  \y{Góc giữa hai vectơ $\vec{a}$ và $\vec{b}$ bằng $90^\circ$.}{d}}

\cauDS[id=C2-B8-D08-NB-046,muc=NB,loigiai={Kiểm tra từng mệnh đề bằng công thức tọa độ. a) Sai. b) Đúng. c) Đúng. d) Sai.}]{Trong không gian $Oxyz$, cho $\vec{a}=(-2;-1;1)$, $\vec{b}=(-1;-2;-1)$. Xét tính đúng sai của các mệnh đề sau:}{%
  \y{Độ dài của vectơ $\vec{a}$ bằng $\sqrt{10}$.}{s}
  \y{Góc giữa hai vectơ $\vec{a}$ và $\vec{b}$ bằng $60^\circ$.}{d}
  \y{$\vec{a}\cdot\vec{b}=3$.}{d}
  \y{Góc giữa hai vectơ $\vec{a}$ và $\vec{b}$ bằng $30^\circ$.}{s}}

\cauDS[id=C2-B8-D08-NB-047,muc=NB,loigiai={Kiểm tra từng mệnh đề bằng công thức tọa độ. a) Sai. b) Sai. c) Đúng. d) Sai.}]{Trong không gian $Oxyz$, cho $\vec{a}=(0;3;2)$, $\vec{b}=(0;-2;3)$. Xét tính đúng sai của các mệnh đề sau:}{%
  \y{Góc giữa hai vectơ $\vec{a}$ và $\vec{b}$ bằng $60^\circ$.}{s}
  \y{Góc giữa hai vectơ $\vec{a}$ và $\vec{b}$ bằng $150^\circ$.}{s}
  \y{Độ dài của vectơ $\vec{a}$ bằng $\sqrt{13}$.}{d}
  \y{$\vec{a}\cdot\vec{b}=-12$.}{s}}

\cauDS[id=C2-B8-D08-NB-048,muc=NB,loigiai={Kiểm tra từng mệnh đề bằng công thức tọa độ. a) Đúng. b) Đúng. c) Đúng. d) Sai.}]{Trong không gian $Oxyz$, cho $\vec{a}=(-3;0;3)$, $\vec{b}=(2;2;-1)$. Xét tính đúng sai của các mệnh đề sau:}{%
  \y{Độ dài của vectơ $\vec{a}$ bằng $3\sqrt{2}$.}{d}
  \y{$\vec{a}\cdot\vec{b}=-9$.}{d}
  \y{Góc giữa hai vectơ $\vec{a}$ và $\vec{b}$ bằng $135^\circ$.}{d}
  \y{Góc giữa hai vectơ $\vec{a}$ và $\vec{b}$ bằng $30^\circ$.}{s}}

\cauTLN[id=C2-B8-D08-NB-049,muc=NB,loigiai={Ta có $\vec{a}\cdot\vec{b}=2\cdot-1+0\cdot2+0\cdot-1=-2$.}]{Trong không gian $Oxyz$, cho $\vec{a}=(2;0;0)$, $\vec{b}=(-1;2;-1)$. Tính $\vec{a}\cdot\vec{b}$.}{-2}

\cauTLN[id=C2-B8-D08-NB-050,muc=NB,loigiai={Ta có $\vec{a}\cdot\vec{b}=1\cdot3-2\cdot-3+1\cdot1=10$.}]{Trong không gian $Oxyz$, cho $\vec{a}=(1;-2;1)$, $\vec{b}=(3;-3;1)$. Tính $\vec{a}\cdot\vec{b}$.}{10}

\cauTLN[id=C2-B8-D08-NB-051,muc=NB,loigiai={Ta có $\vec{a}\cdot\vec{b}=3\cdot3-2\cdot0-1\cdot-2=11$.}]{Trong không gian $Oxyz$, cho $\vec{a}=(3;-2;-1)$, $\vec{b}=(3;0;-2)$. Tính $\vec{a}\cdot\vec{b}$.}{11}

\cauTLN[id=C2-B8-D08-NB-052,muc=NB,loigiai={Ta có $\vec{a}\cdot\vec{b}=-2\cdot0-2\cdot0+1\cdot2=2$.}]{Trong không gian $Oxyz$, cho $\vec{a}=(-2;-2;1)$, $\vec{b}=(0;0;2)$. Tính $\vec{a}\cdot\vec{b}$.}{2}

\cauTLN[id=C2-B8-D08-NB-053,muc=NB,loigiai={Ta có $\vec{a}\cdot\vec{b}=-1\cdot1-1\cdot0+0\cdot3=-1$.}]{Trong không gian $Oxyz$, cho $\vec{a}=(-1;-1;0)$, $\vec{b}=(1;0;3)$. Tính $\vec{a}\cdot\vec{b}$.}{-1}

\cauTLN[id=C2-B8-D08-NB-054,muc=NB,loigiai={Ta có $\vec{a}\cdot\vec{b}=1\cdot-3+3\cdot-1+1\cdot-1=-7$.}]{Trong không gian $Oxyz$, cho $\vec{a}=(1;3;1)$, $\vec{b}=(-3;-1;-1)$. Tính $\vec{a}\cdot\vec{b}$.}{-7}

\cauTLN[id=C2-B8-D08-NB-055,muc=NB,loigiai={Ta có $\cos(\vec{a},\vec{b})=\dfrac{\vec{a}\cdot\vec{b}}{|\vec{a}|\,|\vec{b}|}=\dfrac{0}{\sqrt{12}\cdot\sqrt{6}}=0$, suy ra góc bằng $90^\circ$.}]{Trong không gian $Oxyz$, cho $\vec{a}=(-2;-2;2)$, $\vec{b}=(2;-1;1)$. Tính số đo (độ) góc giữa hai vectơ $\vec{a}$ và $\vec{b}$.}{90}

\cauTLN[id=C2-B8-D08-NB-056,muc=NB,loigiai={Ta có $\cos(\vec{a},\vec{b})=\dfrac{\vec{a}\cdot\vec{b}}{|\vec{a}|\,|\vec{b}|}=\dfrac{0}{\sqrt{2}\cdot\sqrt{11}}=0$, suy ra góc bằng $90^\circ$.}]{Trong không gian $Oxyz$, cho $\vec{a}=(-1;-1;0)$, $\vec{b}=(-1;1;-3)$. Tính số đo (độ) góc giữa hai vectơ $\vec{a}$ và $\vec{b}$.}{90}

\cauTLN[id=C2-B8-D08-NB-057,muc=NB,loigiai={Ta có $\cos(\vec{a},\vec{b})=\dfrac{\vec{a}\cdot\vec{b}}{|\vec{a}|\,|\vec{b}|}=\dfrac{0}{\sqrt{3}\cdot\sqrt{14}}=0$, suy ra góc bằng $90^\circ$.}]{Trong không gian $Oxyz$, cho $\vec{a}=(-1;-1;-1)$, $\vec{b}=(1;2;-3)$. Tính số đo (độ) góc giữa hai vectơ $\vec{a}$ và $\vec{b}$.}{90}

\cauTLN[id=C2-B8-D08-NB-058,muc=NB,loigiai={Ta có $\cos(\vec{a},\vec{b})=\dfrac{\vec{a}\cdot\vec{b}}{|\vec{a}|\,|\vec{b}|}=\dfrac{0}{\sqrt{10}\cdot\sqrt{14}}=0$, suy ra góc bằng $90^\circ$.}]{Trong không gian $Oxyz$, cho $\vec{a}=(0;3;1)$, $\vec{b}=(-2;-1;3)$. Tính số đo (độ) góc giữa hai vectơ $\vec{a}$ và $\vec{b}$.}{90}

\cauTLN[id=C2-B8-D08-NB-059,muc=NB,loigiai={Ta có $\cos(\vec{a},\vec{b})=\dfrac{\vec{a}\cdot\vec{b}}{|\vec{a}|\,|\vec{b}|}=\dfrac{-5}{\sqrt{5}\cdot\sqrt{10}}=- \dfrac{\sqrt{2}}{2}$, suy ra góc bằng $135^\circ$.}]{Trong không gian $Oxyz$, cho $\vec{a}=(0;1;2)$, $\vec{b}=(0;1;-3)$. Tính số đo (độ) góc giữa hai vectơ $\vec{a}$ và $\vec{b}$.}{135}

\cauTLN[id=C2-B8-D08-NB-060,muc=NB,loigiai={Ta có $\cos(\vec{a},\vec{b})=\dfrac{\vec{a}\cdot\vec{b}}{|\vec{a}|\,|\vec{b}|}=\dfrac{-3}{\sqrt{1}\cdot\sqrt{18}}=- \dfrac{\sqrt{2}}{2}$, suy ra góc bằng $135^\circ$.}]{Trong không gian $Oxyz$, cho $\vec{a}=(0;0;1)$, $\vec{b}=(0;3;-3)$. Tính số đo (độ) góc giữa hai vectơ $\vec{a}$ và $\vec{b}$.}{135}

\cauTN[id=C2-B8-D08-TH-001,muc=TH,loigiai={Ta có $\vec{a}\cdot\vec{b}=-14$, $|\vec{a}|=\sqrt{66}$, $|\vec{b}|=\sqrt{35}$ nên $\cos(\vec{a},\vec{b})=\dfrac{-14}{\sqrt{66}\sqrt{35}}=-\dfrac{\sqrt{2310}}{165}$.}]{Cho $\vec{a}=(5;5;-4)$, $\vec{b}=(-5;3;1)$. Giá trị $\cos(\vec{a},\vec{b})$ bằng}
  {$-\dfrac{\sqrt{2310}}{165}$.}
  {$-\dfrac{13\sqrt{2310}}{2310}$.}{$\dfrac{\sqrt{2310}}{165}$.}{$-\dfrac{\sqrt{2310}}{154}$.}

\cauTN[id=C2-B8-D08-TH-002,muc=TH,loigiai={Ta có $\vec{a}\cdot\vec{b}=20$, $|\vec{a}|=\sqrt{35}$, $|\vec{b}|=\sqrt{30}$ nên $\cos(\vec{a},\vec{b})=\dfrac{20}{\sqrt{35}\sqrt{30}}=\dfrac{2\sqrt{42}}{21}$.}]{Cho $\vec{a}=(-3;5;-1)$, $\vec{b}=(1;5;2)$. Giá trị $\cos(\vec{a},\vec{b})$ bằng}
  {$\dfrac{2\sqrt{42}}{21}$.}
  {$\dfrac{\sqrt{42}}{10}$.}{$-\dfrac{2\sqrt{42}}{21}$.}{$\dfrac{19\sqrt{42}}{210}$.}

\cauTN[id=C2-B8-D08-TH-003,muc=TH,loigiai={Ta có $\vec{a}\cdot\vec{b}=17$, $|\vec{a}|=\sqrt{29}$, $|\vec{b}|=\sqrt{13}$ nên $\cos(\vec{a},\vec{b})=\dfrac{17}{\sqrt{29}\sqrt{13}}=\dfrac{17\sqrt{377}}{377}$.}]{Cho $\vec{a}=(2;-4;3)$, $\vec{b}=(0;-2;3)$. Giá trị $\cos(\vec{a},\vec{b})$ bằng}
  {$\dfrac{17\sqrt{377}}{377}$.}
  {$\dfrac{18\sqrt{377}}{377}$.}{$-\dfrac{17\sqrt{377}}{377}$.}{$\dfrac{16\sqrt{377}}{377}$.}

\cauTN[id=C2-B8-D08-TH-004,muc=TH,loigiai={Ta có $\vec{a}\cdot\vec{b}=18$, $|\vec{a}|=\sqrt{8}$, $|\vec{b}|=\sqrt{45}$ nên $\cos(\vec{a},\vec{b})=\dfrac{18}{\sqrt{8}\sqrt{45}}=\dfrac{3\sqrt{10}}{10}$.}]{Cho $\vec{a}=(2;2;0)$, $\vec{b}=(4;5;2)$. Giá trị $\cos(\vec{a},\vec{b})$ bằng}
  {$\dfrac{3\sqrt{10}}{10}$.}
  {$-\dfrac{3\sqrt{10}}{10}$.}{$\dfrac{17\sqrt{10}}{60}$.}{$\dfrac{2\sqrt{5}}{5}$.}

\cauTN[id=C2-B8-D08-TH-005,muc=TH,loigiai={Ta có $\vec{a}\cdot\vec{b}=-15$, $|\vec{a}|=\sqrt{45}$, $|\vec{b}|=\sqrt{35}$ nên $\cos(\vec{a},\vec{b})=\dfrac{-15}{\sqrt{45}\sqrt{35}}=-\dfrac{\sqrt{7}}{7}$.}]{Cho $\vec{a}=(2;5;-4)$, $\vec{b}=(3;-5;-1)$. Giá trị $\cos(\vec{a},\vec{b})$ bằng}
  {$-\dfrac{\sqrt{7}}{7}$.}
  {$-\dfrac{2\sqrt{7}}{15}$.}{$\dfrac{\sqrt{7}}{7}$.}{$-\dfrac{16\sqrt{7}}{105}$.}

\cauTN[id=C2-B8-D08-TH-006,muc=TH,loigiai={Ta có $\vec{a}\cdot\vec{b}=-1$, $|\vec{a}|=\sqrt{14}$, $|\vec{b}|=\sqrt{29}$ nên $\cos(\vec{a},\vec{b})=\dfrac{-1}{\sqrt{14}\sqrt{29}}=-\dfrac{\sqrt{406}}{406}$.}]{Cho $\vec{a}=(1;-3;-2)$, $\vec{b}=(-5;0;-2)$. Giá trị $\cos(\vec{a},\vec{b})$ bằng}
  {$-\dfrac{\sqrt{406}}{406}$.}
  {$0$.}{$\dfrac{\sqrt{406}}{406}$.}{$-\dfrac{\sqrt{406}}{203}$.}

\cauTN[id=C2-B8-D08-TH-007,muc=TH,loigiai={Ta có $\vec{a}\cdot\vec{b}=-14$, $|\vec{a}|=\sqrt{27}$, $|\vec{b}|=\sqrt{22}$ nên $\cos(\vec{a},\vec{b})=\dfrac{-14}{\sqrt{27}\sqrt{22}}=-\dfrac{7\sqrt{66}}{99}$.}]{Cho $\vec{a}=(1;5;1)$, $\vec{b}=(3;-3;-2)$. Giá trị $\cos(\vec{a},\vec{b})$ bằng}
  {$-\dfrac{7\sqrt{66}}{99}$.}
  {$-\dfrac{13\sqrt{66}}{198}$.}{$\dfrac{7\sqrt{66}}{99}$.}{$-\dfrac{5\sqrt{66}}{66}$.}

\cauTN[id=C2-B8-D08-TH-008,muc=TH,loigiai={Ta có $\vec{a}\cdot\vec{b}=12$, $|\vec{a}|=\sqrt{20}$, $|\vec{b}|=\sqrt{50}$ nên $\cos(\vec{a},\vec{b})=\dfrac{12}{\sqrt{20}\sqrt{50}}=\dfrac{3\sqrt{10}}{25}$.}]{Cho $\vec{a}=(-4;0;-2)$, $\vec{b}=(-5;3;4)$. Giá trị $\cos(\vec{a},\vec{b})$ bằng}
  {$\dfrac{3\sqrt{10}}{25}$.}
  {$\dfrac{13\sqrt{10}}{100}$.}{$-\dfrac{3\sqrt{10}}{25}$.}{$\dfrac{11\sqrt{10}}{100}$.}

\cauTN[id=C2-B8-D08-TH-009,muc=TH,loigiai={Ta có $\vec{a}\cdot\vec{b}=-3$, $|\vec{a}|=\sqrt{26}$, $|\vec{b}|=\sqrt{29}$ nên $\cos(\vec{a},\vec{b})=\dfrac{-3}{\sqrt{26}\sqrt{29}}=-\dfrac{3\sqrt{754}}{754}$.}]{Cho $\vec{a}=(-3;4;1)$, $\vec{b}=(-3;-2;-4)$. Giá trị $\cos(\vec{a},\vec{b})$ bằng}
  {$-\dfrac{3\sqrt{754}}{754}$.}
  {$-\dfrac{\sqrt{754}}{377}$.}{$\dfrac{3\sqrt{754}}{754}$.}{$-\dfrac{2\sqrt{754}}{377}$.}

\cauTN[id=C2-B8-D08-TH-010,muc=TH,loigiai={Ta có $\vec{a}\cdot\vec{b}=-34$, $|\vec{a}|=\sqrt{51}$, $|\vec{b}|=\sqrt{34}$ nên $\cos(\vec{a},\vec{b})=\dfrac{-34}{\sqrt{51}\sqrt{34}}=-\dfrac{\sqrt{6}}{3}$.}]{Cho $\vec{a}=(-1;5;5)$, $\vec{b}=(4;-3;-3)$. Giá trị $\cos(\vec{a},\vec{b})$ bằng}
  {$-\dfrac{\sqrt{6}}{3}$.}
  {$-\dfrac{11\sqrt{6}}{34}$.}{$\dfrac{\sqrt{6}}{3}$.}{$-\dfrac{35\sqrt{6}}{102}$.}

\cauTN[id=C2-B8-D08-TH-011,muc=TH,loigiai={Ta có $\vec{a}\cdot\vec{b}=3$, $|\vec{a}|=\sqrt{42}$, $|\vec{b}|=\sqrt{14}$ nên $\cos(\vec{a},\vec{b})=\dfrac{3}{\sqrt{42}\sqrt{14}}=\dfrac{\sqrt{3}}{14}$.}]{Cho $\vec{a}=(5;-1;4)$, $\vec{b}=(2;3;-1)$. Giá trị $\cos(\vec{a},\vec{b})$ bằng}
  {$\dfrac{\sqrt{3}}{14}$.}
  {$\dfrac{2\sqrt{3}}{21}$.}{$-\dfrac{\sqrt{3}}{14}$.}{$\dfrac{\sqrt{3}}{21}$.}

\cauTN[id=C2-B8-D08-TH-012,muc=TH,loigiai={Ta có $\vec{a}\cdot\vec{b}=6$, $|\vec{a}|=\sqrt{11}$, $|\vec{b}|=\sqrt{22}$ nên $\cos(\vec{a},\vec{b})=\dfrac{6}{\sqrt{11}\sqrt{22}}=\dfrac{3\sqrt{2}}{11}$.}]{Cho $\vec{a}=(-1;-3;1)$, $\vec{b}=(-3;-2;-3)$. Giá trị $\cos(\vec{a},\vec{b})$ bằng}
  {$\dfrac{3\sqrt{2}}{11}$.}
  {$\dfrac{7\sqrt{2}}{22}$.}{$-\dfrac{3\sqrt{2}}{11}$.}{$\dfrac{5\sqrt{2}}{22}$.}

\cauTN[id=C2-B8-D08-TH-013,muc=TH,loigiai={Ta có $3\vec{a}+\vec{b}=(15;1;9)$, $\vec{a}-\vec{b}=(5;-5;3)$ nên tích vô hướng bằng $97$.}]{Cho $\vec{a}=(5;-1;3)$, $\vec{b}=(0;4;0)$. Tích vô hướng $(3\vec{a}+\vec{b})\cdot(\vec{a}-\vec{b})$ bằng}
  {$97$.}
  {$\dfrac{97}{2}$.}{$195$.}{$95$.}

\cauTN[id=C2-B8-D08-TH-014,muc=TH,loigiai={Ta có $\vec{a}+\vec{b}=(-2;-8;-8)$, $\vec{a}+2\vec{b}=(-7;-11;-12)$ nên tích vô hướng bằng $198$.}]{Cho $\vec{a}=(3;-5;-4)$, $\vec{b}=(-5;-3;-4)$. Tích vô hướng $(\vec{a}+\vec{b})\cdot(\vec{a}+2\vec{b})$ bằng}
  {$198$.}
  {$197$.}{$99$.}{$397$.}

\cauTN[id=C2-B8-D08-TH-015,muc=TH,loigiai={Ta có $-\vec{a}+\vec{b}=(-1;-6;-2)$, $3\vec{a}+2\vec{b}=(3;3;-14)$ nên tích vô hướng bằng $7$.}]{Cho $\vec{a}=(1;3;-2)$, $\vec{b}=(0;-3;-4)$. Tích vô hướng $(-\vec{a}+\vec{b})\cdot(3\vec{a}+2\vec{b})$ bằng}
  {$7$.}
  {$\dfrac{7}{2}$.}{$15$.}{$5$.}

\cauTN[id=C2-B8-D08-TH-016,muc=TH,loigiai={Ta có $2\vec{a}+\vec{b}=(-7;1;7)$, $2\vec{a}+3\vec{b}=(-1;3;1)$ nên tích vô hướng bằng $17$.}]{Cho $\vec{a}=(-5;0;5)$, $\vec{b}=(3;1;-3)$. Tích vô hướng $(2\vec{a}+\vec{b})\cdot(2\vec{a}+3\vec{b})$ bằng}
  {$17$.}
  {$16$.}{$35$.}{$15$.}

\cauTN[id=C2-B8-D08-TH-017,muc=TH,loigiai={Ta có $2\vec{a}+\vec{b}=(1;-7;13)$, $2\vec{a}+\vec{b}=(1;-7;13)$ nên tích vô hướng bằng $219$.}]{Cho $\vec{a}=(-1;-1;5)$, $\vec{b}=(3;-5;3)$. Tích vô hướng $(2\vec{a}+\vec{b})\cdot(2\vec{a}+\vec{b})$ bằng}
  {$219$.}
  {$439$.}{$220$.}{$221$.}

\cauTN[id=C2-B8-D08-TH-018,muc=TH,loigiai={Ta có $2\vec{a}+2\vec{b}=(4;8;-10)$, $\vec{a}+\vec{b}=(2;4;-5)$ nên tích vô hướng bằng $90$.}]{Cho $\vec{a}=(5;4;-1)$, $\vec{b}=(-3;0;-4)$. Tích vô hướng $(2\vec{a}+2\vec{b})\cdot(\vec{a}+\vec{b})$ bằng}
  {$90$.}
  {$89$.}{$45$.}{$180$.}

\cauTN[id=C2-B8-D08-TH-019,muc=TH,loigiai={Ta có $2\vec{a}-\vec{b}=(-10;11;-2)$, $\vec{a}-\vec{b}=(-5;7;0)$ nên tích vô hướng bằng $127$.}]{Cho $\vec{a}=(-5;4;-2)$, $\vec{b}=(0;-3;-2)$. Tích vô hướng $(2\vec{a}-\vec{b})\cdot(\vec{a}-\vec{b})$ bằng}
  {$127$.}
  {$124$.}{$\dfrac{127}{2}$.}{$129$.}

\cauTN[id=C2-B8-D08-TH-020,muc=TH,loigiai={Ta có $\vec{a}+3\vec{b}=(3;3;2)$, $\vec{a}-\vec{b}=(-1;-1;-2)$ nên tích vô hướng bằng $-10$.}]{Cho $\vec{a}=(0;0;-1)$, $\vec{b}=(1;1;1)$. Tích vô hướng $(\vec{a}+3\vec{b})\cdot(\vec{a}-\vec{b})$ bằng}
  {$-10$.}
  {$-12$.}{$-19$.}{$-7$.}

\cauTN[id=C2-B8-D08-TH-021,muc=TH,loigiai={Ta có $-\vec{a}+\vec{b}=(6;4;-4)$, $\vec{a}+3\vec{b}=(14;-8;4)$ nên tích vô hướng bằng $36$.}]{Cho $\vec{a}=(-1;-5;4)$, $\vec{b}=(5;-1;0)$. Tích vô hướng $(-\vec{a}+\vec{b})\cdot(\vec{a}+3\vec{b})$ bằng}
  {$36$.}
  {$34$.}{$38$.}{$35$.}

\cauTN[id=C2-B8-D08-TH-022,muc=TH,loigiai={Ta có $-\vec{a}-\vec{b}=(-2;1;-1)$, $\vec{a}+2\vec{b}=(5;3;0)$ nên tích vô hướng bằng $-7$.}]{Cho $\vec{a}=(-1;-5;2)$, $\vec{b}=(3;4;-1)$. Tích vô hướng $(-\vec{a}-\vec{b})\cdot(\vec{a}+2\vec{b})$ bằng}
  {$-7$.}
  {$-5$.}{$-\dfrac{7}{2}$.}{$-9$.}

\cauTN[id=C2-B8-D08-TH-023,muc=TH,loigiai={Ta có $\vec{a}+\vec{b}=(-8;-8;-1)$, $\vec{a}+3\vec{b}=(-16;-18;-3)$ nên tích vô hướng bằng $275$.}]{Cho $\vec{a}=(-4;-3;0)$, $\vec{b}=(-4;-5;-1)$. Tích vô hướng $(\vec{a}+\vec{b})\cdot(\vec{a}+3\vec{b})$ bằng}
  {$275$.}
  {$550$.}{$274$.}{$272$.}

\cauTN[id=C2-B8-D08-TH-024,muc=TH,loigiai={Ta có $-\vec{a}+2\vec{b}=(7;7;5)$, $2\vec{a}-\vec{b}=(-2;-8;-1)$ nên tích vô hướng bằng $-75$.}]{Cho $\vec{a}=(1;-3;1)$, $\vec{b}=(4;2;3)$. Tích vô hướng $(-\vec{a}+2\vec{b})\cdot(2\vec{a}-\vec{b})$ bằng}
  {$-75$.}
  {$-\dfrac{75}{2}$.}{$-72$.}{$-73$.}

\cauTN[id=C2-B8-D08-TH-025,muc=TH,loigiai={Hình chiếu của $\vec{a}$ lên phương của $\vec{b}$ là $\dfrac{\vec{a}\cdot\vec{b}}{|\vec{b}|^2}\vec{b}=\dfrac{-18}{34}\vec{b}=(0;-\dfrac{45}{17};\dfrac{27}{17})$.}]{Cho $\vec{a}=(4;-3;1)$, $\vec{b}=(0;5;-3)$. Hình chiếu vuông góc của $\vec{a}$ lên phương của $\vec{b}$ là vectơ}
  {$(0;-\dfrac{45}{17};\dfrac{27}{17})$.}
  {$(0;-\dfrac{62}{17};\dfrac{27}{17})$.}{$(0;-\dfrac{11}{17};\dfrac{27}{17})$.}{$(0;\dfrac{27}{17};-\dfrac{45}{17})$.}

\cauTN[id=C2-B8-D08-TH-026,muc=TH,loigiai={Hình chiếu của $\vec{a}$ lên phương của $\vec{b}$ là $\dfrac{\vec{a}\cdot\vec{b}}{|\vec{b}|^2}\vec{b}=\dfrac{47}{51}\vec{b}=(-\dfrac{235}{51};-\dfrac{235}{51};-\dfrac{47}{51})$.}]{Cho $\vec{a}=(-5;-4;-2)$, $\vec{b}=(-5;-5;-1)$. Hình chiếu vuông góc của $\vec{a}$ lên phương của $\vec{b}$ là vectơ}
  {$(-\dfrac{235}{51};-\dfrac{235}{51};-\dfrac{47}{51})$.}
  {$(-\dfrac{235}{51};\dfrac{235}{51};-\dfrac{47}{51})$.}{$(-\dfrac{184}{51};-\dfrac{235}{51};-\dfrac{47}{51})$.}{$(-\dfrac{235}{51};-\dfrac{235}{51};\dfrac{4}{51})$.}

\cauTN[id=C2-B8-D08-TH-027,muc=TH,loigiai={Hình chiếu của $\vec{a}$ lên phương của $\vec{b}$ là $\dfrac{\vec{a}\cdot\vec{b}}{|\vec{b}|^2}\vec{b}=\dfrac{36}{59}\vec{b}=(-\dfrac{180}{59};-\dfrac{180}{59};\dfrac{108}{59})$.}]{Cho $\vec{a}=(-4;-5;-3)$, $\vec{b}=(-5;-5;3)$. Hình chiếu vuông góc của $\vec{a}$ lên phương của $\vec{b}$ là vectơ}
  {$(-\dfrac{180}{59};-\dfrac{180}{59};\dfrac{108}{59})$.}
  {$(-\dfrac{180}{59};-\dfrac{298}{59};\dfrac{108}{59})$.}{$(-\dfrac{180}{59};-\dfrac{62}{59};\dfrac{108}{59})$.}{$(-\dfrac{180}{59};-\dfrac{180}{59};\dfrac{49}{59})$.}

\cauTN[id=C2-B8-D08-TH-028,muc=TH,loigiai={Hình chiếu của $\vec{a}$ lên phương của $\vec{b}$ là $\dfrac{\vec{a}\cdot\vec{b}}{|\vec{b}|^2}\vec{b}=\dfrac{-21}{30}\vec{b}=(\dfrac{7}{2};\dfrac{7}{5};-\dfrac{7}{10})$.}]{Cho $\vec{a}=(5;-4;-4)$, $\vec{b}=(-5;-2;1)$. Hình chiếu vuông góc của $\vec{a}$ lên phương của $\vec{b}$ là vectơ}
  {$(\dfrac{7}{2};\dfrac{7}{5};-\dfrac{7}{10})$.}
  {$(-\dfrac{7}{10};\dfrac{7}{5};\dfrac{7}{2})$.}{$(-\dfrac{7}{2};-\dfrac{7}{5};\dfrac{7}{10})$.}{$(\dfrac{7}{2};\dfrac{12}{5};-\dfrac{7}{10})$.}

\cauTN[id=C2-B8-D08-TH-029,muc=TH,loigiai={Hình chiếu của $\vec{a}$ lên phương của $\vec{b}$ là $\dfrac{\vec{a}\cdot\vec{b}}{|\vec{b}|^2}\vec{b}=\dfrac{5}{22}\vec{b}=(\dfrac{15}{22};\dfrac{15}{22};-\dfrac{5}{11})$.}]{Cho $\vec{a}=(2;-1;-1)$, $\vec{b}=(3;3;-2)$. Hình chiếu vuông góc của $\vec{a}$ lên phương của $\vec{b}$ là vectơ}
  {$(\dfrac{15}{22};\dfrac{15}{22};-\dfrac{5}{11})$.}
  {$(\dfrac{15}{22};\dfrac{15}{22};-\dfrac{27}{11})$.}{$(-\dfrac{15}{22};-\dfrac{15}{22};\dfrac{5}{11})$.}{$(15;15;-10)$.}

\cauTN[id=C2-B8-D08-TH-030,muc=TH,loigiai={Hình chiếu của $\vec{a}$ lên phương của $\vec{b}$ là $\dfrac{\vec{a}\cdot\vec{b}}{|\vec{b}|^2}\vec{b}=\dfrac{6}{27}\vec{b}=(-\dfrac{2}{9};\dfrac{10}{9};\dfrac{2}{9})$.}]{Cho $\vec{a}=(-3;0;3)$, $\vec{b}=(-1;5;1)$. Hình chiếu vuông góc của $\vec{a}$ lên phương của $\vec{b}$ là vectơ}
  {$(-\dfrac{2}{9};\dfrac{10}{9};\dfrac{2}{9})$.}
  {$(-\dfrac{2}{9};-\dfrac{10}{9};\dfrac{2}{9})$.}{$(-\dfrac{2}{9};\dfrac{28}{9};\dfrac{2}{9})$.}{$(\dfrac{10}{9};-\dfrac{2}{9};\dfrac{2}{9})$.}

\cauTN[id=C2-B8-D08-TH-031,muc=TH,loigiai={Hình chiếu của $\vec{a}$ lên phương của $\vec{b}$ là $\dfrac{\vec{a}\cdot\vec{b}}{|\vec{b}|^2}\vec{b}=\dfrac{-14}{9}\vec{b}=(\dfrac{28}{9};\dfrac{14}{9};-\dfrac{28}{9})$.}]{Cho $\vec{a}=(4;4;-1)$, $\vec{b}=(-2;-1;2)$. Hình chiếu vuông góc của $\vec{a}$ lên phương của $\vec{b}$ là vectơ}
  {$(\dfrac{28}{9};\dfrac{14}{9};-\dfrac{28}{9})$.}
  {$(\dfrac{46}{9};\dfrac{14}{9};-\dfrac{28}{9})$.}{$(\dfrac{28}{9};\dfrac{14}{9};\dfrac{28}{9})$.}{$(\dfrac{28}{9};-\dfrac{14}{9};-\dfrac{28}{9})$.}

\cauTN[id=C2-B8-D08-TH-032,muc=TH,loigiai={Hình chiếu của $\vec{a}$ lên phương của $\vec{b}$ là $\dfrac{\vec{a}\cdot\vec{b}}{|\vec{b}|^2}\vec{b}=\dfrac{4}{30}\vec{b}=(-\dfrac{2}{3};\dfrac{2}{15};\dfrac{4}{15})$.}]{Cho $\vec{a}=(0;2;1)$, $\vec{b}=(-5;1;2)$. Hình chiếu vuông góc của $\vec{a}$ lên phương của $\vec{b}$ là vectơ}
  {$(-\dfrac{2}{3};\dfrac{2}{15};\dfrac{4}{15})$.}
  {$(-\dfrac{2}{3};-\dfrac{13}{15};\dfrac{4}{15})$.}{$(\dfrac{4}{15};\dfrac{2}{15};-\dfrac{2}{3})$.}{$(-\dfrac{2}{3};\dfrac{2}{15};\dfrac{34}{15})$.}

\cauTN[id=C2-B8-D08-TH-033,muc=TH,loigiai={Hình chiếu của $\vec{a}$ lên phương của $\vec{b}$ là $\dfrac{\vec{a}\cdot\vec{b}}{|\vec{b}|^2}\vec{b}=\dfrac{17}{30}\vec{b}=(\dfrac{17}{15};-\dfrac{17}{30};-\dfrac{17}{6})$.}]{Cho $\vec{a}=(3;-1;-2)$, $\vec{b}=(2;-1;-5)$. Hình chiếu vuông góc của $\vec{a}$ lên phương của $\vec{b}$ là vectơ}
  {$(\dfrac{17}{15};-\dfrac{17}{30};-\dfrac{17}{6})$.}
  {$(\dfrac{17}{15};-\dfrac{17}{30};-\dfrac{11}{6})$.}{$(-\dfrac{17}{15};\dfrac{17}{30};\dfrac{17}{6})$.}{$(-\dfrac{17}{30};\dfrac{17}{15};-\dfrac{17}{6})$.}

\cauTN[id=C2-B8-D08-TH-034,muc=TH,loigiai={Hình chiếu của $\vec{a}$ lên phương của $\vec{b}$ là $\dfrac{\vec{a}\cdot\vec{b}}{|\vec{b}|^2}\vec{b}=\dfrac{26}{19}\vec{b}=(\dfrac{26}{19};\dfrac{78}{19};\dfrac{78}{19})$.}]{Cho $\vec{a}=(-1;4;5)$, $\vec{b}=(1;3;3)$. Hình chiếu vuông góc của $\vec{a}$ lên phương của $\vec{b}$ là vectơ}
  {$(\dfrac{26}{19};\dfrac{78}{19};\dfrac{78}{19})$.}
  {$(\dfrac{7}{19};\dfrac{78}{19};\dfrac{78}{19})$.}{$(\dfrac{26}{19};-\dfrac{78}{19};\dfrac{78}{19})$.}{$(\dfrac{26}{19};\dfrac{78}{19};\dfrac{97}{19})$.}

\cauTN[id=C2-B8-D08-TH-035,muc=TH,loigiai={Hình chiếu của $\vec{a}$ lên phương của $\vec{b}$ là $\dfrac{\vec{a}\cdot\vec{b}}{|\vec{b}|^2}\vec{b}=\dfrac{-7}{19}\vec{b}=(-\dfrac{21}{19};\dfrac{7}{19};-\dfrac{21}{19})$.}]{Cho $\vec{a}=(-2;1;0)$, $\vec{b}=(3;-1;3)$. Hình chiếu vuông góc của $\vec{a}$ lên phương của $\vec{b}$ là vectơ}
  {$(-\dfrac{21}{19};\dfrac{7}{19};-\dfrac{21}{19})$.}
  {$(-\dfrac{21}{19};-\dfrac{31}{19};-\dfrac{21}{19})$.}{$(\dfrac{7}{19};-\dfrac{21}{19};-\dfrac{21}{19})$.}{$(-\dfrac{21}{19};-\dfrac{12}{19};-\dfrac{21}{19})$.}

\cauTN[id=C2-B8-D08-TH-036,muc=TH,loigiai={Hình chiếu của $\vec{a}$ lên phương của $\vec{b}$ là $\dfrac{\vec{a}\cdot\vec{b}}{|\vec{b}|^2}\vec{b}=\dfrac{-20}{27}\vec{b}=(\dfrac{20}{27};\dfrac{100}{27};\dfrac{20}{27})$.}]{Cho $\vec{a}=(-1;4;1)$, $\vec{b}=(-1;-5;-1)$. Hình chiếu vuông góc của $\vec{a}$ lên phương của $\vec{b}$ là vectơ}
  {$(\dfrac{20}{27};\dfrac{100}{27};\dfrac{20}{27})$.}
  {$(\dfrac{20}{27};\dfrac{100}{27};\dfrac{74}{27})$.}{$(-\dfrac{34}{27};\dfrac{100}{27};\dfrac{20}{27})$.}{$(\dfrac{20}{27};\dfrac{100}{27};-\dfrac{7}{27})$.}

\cauTN[id=C2-B8-D08-TH-037,muc=TH,loigiai={Ta có $\overrightarrow{BA}=(1;-2;9)$, $\overrightarrow{BC}=(-1;-10;6)$ nên $\cos\widehat{ABC}=\dfrac{\overrightarrow{BA}\cdot\overrightarrow{BC}}{|\overrightarrow{BA}||\overrightarrow{BC}|}=\dfrac{73\sqrt{11782}}{11782}$.}]{Cho tam giác $ABC$ với $A(3;3;4)$, $B(2;5;-5)$, $C(1;-5;1)$. Giá trị $\cos\widehat{ABC}$ bằng}
  {$\dfrac{73\sqrt{11782}}{11782}$.}
  {$\dfrac{37\sqrt{11782}}{5891}$.}{$-\dfrac{73\sqrt{11782}}{11782}$.}{$\dfrac{36\sqrt{11782}}{5891}$.}

\cauTN[id=C2-B8-D08-TH-038,muc=TH,loigiai={Ta có $\overrightarrow{BA}=(3;8;-2)$, $\overrightarrow{BC}=(9;6;3)$ nên $\cos\widehat{ABC}=\dfrac{\overrightarrow{BA}\cdot\overrightarrow{BC}}{|\overrightarrow{BA}||\overrightarrow{BC}|}=\dfrac{23\sqrt{22}}{154}$.}]{Cho tam giác $ABC$ với $A(-1;4;-4)$, $B(-4;-4;-2)$, $C(5;2;1)$. Giá trị $\cos\widehat{ABC}$ bằng}
  {$\dfrac{23\sqrt{22}}{154}$.}
  {$\dfrac{5\sqrt{22}}{33}$.}{$-\dfrac{23\sqrt{22}}{154}$.}{$\dfrac{34\sqrt{22}}{231}$.}

\cauTN[id=C2-B8-D08-TH-039,muc=TH,loigiai={Ta có $\overrightarrow{BA}=(0;-3;3)$, $\overrightarrow{BC}=(-4;-4;-2)$ nên $\cos\widehat{ABC}=\dfrac{\overrightarrow{BA}\cdot\overrightarrow{BC}}{|\overrightarrow{BA}||\overrightarrow{BC}|}=\dfrac{\sqrt{2}}{6}$.}]{Cho tam giác $ABC$ với $A(0;-3;2)$, $B(0;0;-1)$, $C(-4;-4;-3)$. Giá trị $\cos\widehat{ABC}$ bằng}
  {$\dfrac{\sqrt{2}}{6}$.}
  {$\dfrac{7\sqrt{2}}{36}$.}{$-\dfrac{\sqrt{2}}{6}$.}{$\dfrac{5\sqrt{2}}{36}$.}

\cauTN[id=C2-B8-D08-TH-040,muc=TH,loigiai={Ta có $\overrightarrow{BA}=(-1;1;-3)$, $\overrightarrow{BC}=(2;0;-4)$ nên $\cos\widehat{ABC}=\dfrac{\overrightarrow{BA}\cdot\overrightarrow{BC}}{|\overrightarrow{BA}||\overrightarrow{BC}|}=\dfrac{\sqrt{55}}{11}$.}]{Cho tam giác $ABC$ với $A(0;0;-1)$, $B(1;-1;2)$, $C(3;-1;-2)$. Giá trị $\cos\widehat{ABC}$ bằng}
  {$\dfrac{\sqrt{55}}{11}$.}
  {$\dfrac{\sqrt{55}}{10}$.}{$-\dfrac{\sqrt{55}}{11}$.}{$\dfrac{9\sqrt{55}}{110}$.}

\cauTN[id=C2-B8-D08-TH-041,muc=TH,loigiai={Ta có $\overrightarrow{BA}=(-7;6;3)$, $\overrightarrow{BC}=(-5;2;-3)$ nên $\cos\widehat{ABC}=\dfrac{\overrightarrow{BA}\cdot\overrightarrow{BC}}{|\overrightarrow{BA}||\overrightarrow{BC}|}=\dfrac{\sqrt{893}}{47}$.}]{Cho tam giác $ABC$ với $A(-3;3;5)$, $B(4;-3;2)$, $C(-1;-1;-1)$. Giá trị $\cos\widehat{ABC}$ bằng}
  {$\dfrac{\sqrt{893}}{47}$.}
  {$\dfrac{39\sqrt{893}}{1786}$.}{$-\dfrac{\sqrt{893}}{47}$.}{$\dfrac{37\sqrt{893}}{1786}$.}

\cauTN[id=C2-B8-D08-TH-042,muc=TH,loigiai={Ta có $\overrightarrow{BA}=(-1;8;-2)$, $\overrightarrow{BC}=(-1;4;-5)$ nên $\cos\widehat{ABC}=\dfrac{\overrightarrow{BA}\cdot\overrightarrow{BC}}{|\overrightarrow{BA}||\overrightarrow{BC}|}=\dfrac{43\sqrt{322}}{966}$.}]{Cho tam giác $ABC$ với $A(1;5;3)$, $B(2;-3;5)$, $C(1;1;0)$. Giá trị $\cos\widehat{ABC}$ bằng}
  {$\dfrac{43\sqrt{322}}{966}$.}
  {$\dfrac{22\sqrt{322}}{483}$.}{$-\dfrac{43\sqrt{322}}{966}$.}{$\dfrac{\sqrt{322}}{23}$.}

\cauTN[id=C2-B8-D08-TH-043,muc=TH,loigiai={Ta có $\overrightarrow{BA}=(2;-9;3)$, $\overrightarrow{BC}=(0;-8;0)$ nên $\cos\widehat{ABC}=\dfrac{\overrightarrow{BA}\cdot\overrightarrow{BC}}{|\overrightarrow{BA}||\overrightarrow{BC}|}=\dfrac{9\sqrt{94}}{94}$.}]{Cho tam giác $ABC$ với $A(-1;-4;5)$, $B(-3;5;2)$, $C(-3;-3;2)$. Giá trị $\cos\widehat{ABC}$ bằng}
  {$\dfrac{9\sqrt{94}}{94}$.}
  {$\dfrac{73\sqrt{94}}{752}$.}{$-\dfrac{9\sqrt{94}}{94}$.}{$\dfrac{71\sqrt{94}}{752}$.}

\cauTN[id=C2-B8-D08-TH-044,muc=TH,loigiai={Ta có $\overrightarrow{BA}=(1;2;-3)$, $\overrightarrow{BC}=(0;8;-1)$ nên $\cos\widehat{ABC}=\dfrac{\overrightarrow{BA}\cdot\overrightarrow{BC}}{|\overrightarrow{BA}||\overrightarrow{BC}|}=\dfrac{19\sqrt{910}}{910}$.}]{Cho tam giác $ABC$ với $A(-4;-1;-1)$, $B(-5;-3;2)$, $C(-5;5;1)$. Giá trị $\cos\widehat{ABC}$ bằng}
  {$\dfrac{19\sqrt{910}}{910}$.}
  {$\dfrac{2\sqrt{910}}{91}$.}{$-\dfrac{19\sqrt{910}}{910}$.}{$\dfrac{9\sqrt{910}}{455}$.}

\cauTN[id=C2-B8-D08-TH-045,muc=TH,loigiai={Ta có $\overrightarrow{BA}=(9;-2;3)$, $\overrightarrow{BC}=(0;6;-7)$ nên $\cos\widehat{ABC}=\dfrac{\overrightarrow{BA}\cdot\overrightarrow{BC}}{|\overrightarrow{BA}||\overrightarrow{BC}|}=-\dfrac{33\sqrt{7990}}{7990}$.}]{Cho tam giác $ABC$ với $A(5;-4;5)$, $B(-4;-2;2)$, $C(-4;4;-5)$. Giá trị $\cos\widehat{ABC}$ bằng}
  {$-\dfrac{33\sqrt{7990}}{7990}$.}
  {$-\dfrac{16\sqrt{7990}}{3995}$.}{$\dfrac{33\sqrt{7990}}{7990}$.}{$-\dfrac{\sqrt{7990}}{235}$.}

\cauTN[id=C2-B8-D08-TH-046,muc=TH,loigiai={Ta có $\overrightarrow{BA}=(4;5;-5)$, $\overrightarrow{BC}=(-4;5;-4)$ nên $\cos\widehat{ABC}=\dfrac{\overrightarrow{BA}\cdot\overrightarrow{BC}}{|\overrightarrow{BA}||\overrightarrow{BC}|}=\dfrac{29\sqrt{418}}{1254}$.}]{Cho tam giác $ABC$ với $A(4;1;-5)$, $B(0;-4;0)$, $C(-4;1;-4)$. Giá trị $\cos\widehat{ABC}$ bằng}
  {$\dfrac{29\sqrt{418}}{1254}$.}
  {$\dfrac{5\sqrt{418}}{209}$.}{$-\dfrac{29\sqrt{418}}{1254}$.}{$\dfrac{14\sqrt{418}}{627}$.}

\cauTN[id=C2-B8-D08-TH-047,muc=TH,loigiai={Ta có $\overrightarrow{BA}=(-3;0;3)$, $\overrightarrow{BC}=(3;-2;-2)$ nên $\cos\widehat{ABC}=\dfrac{\overrightarrow{BA}\cdot\overrightarrow{BC}}{|\overrightarrow{BA}||\overrightarrow{BC}|}=-\dfrac{5\sqrt{34}}{34}$.}]{Cho tam giác $ABC$ với $A(-2;1;0)$, $B(1;1;-3)$, $C(4;-1;-5)$. Giá trị $\cos\widehat{ABC}$ bằng}
  {$-\dfrac{5\sqrt{34}}{34}$.}
  {$-\dfrac{7\sqrt{34}}{51}$.}{$\dfrac{5\sqrt{34}}{34}$.}{$-\dfrac{8\sqrt{34}}{51}$.}

\cauTN[id=C2-B8-D08-TH-048,muc=TH,loigiai={Ta có $\overrightarrow{BA}=(-6;-3;1)$, $\overrightarrow{BC}=(-2;6;3)$ nên $\cos\widehat{ABC}=\dfrac{\overrightarrow{BA}\cdot\overrightarrow{BC}}{|\overrightarrow{BA}||\overrightarrow{BC}|}=-\dfrac{3\sqrt{46}}{322}$.}]{Cho tam giác $ABC$ với $A(-1;-5;-1)$, $B(5;-2;-2)$, $C(3;4;1)$. Giá trị $\cos\widehat{ABC}$ bằng}
  {$-\dfrac{3\sqrt{46}}{322}$.}
  {$-\dfrac{\sqrt{46}}{161}$.}{$\dfrac{3\sqrt{46}}{322}$.}{$-\dfrac{2\sqrt{46}}{161}$.}

\cauDS[id=C2-B8-D08-TH-049,muc=TH,loigiai={Kiểm tra từng mệnh đề bằng công thức tọa độ. a) Đúng. b) Sai. c) Đúng. d) Sai.}]{Trong không gian $Oxyz$, cho $\vec{a}=(0;0;3)$, $\vec{b}=(-1;-4;-5)$. Xét tính đúng sai của các mệnh đề sau:}{%
  \y{Với $\vec{a}=(0;0;3)$, $\vec{b}=(-1;-4;-5)$ thì $(3\vec{a}-\vec{b})\cdot(3\vec{a}-\vec{b})=213$.}{d}
  \y{Với $\vec{a}=(0;0;3)$, $\vec{b}=(-1;-4;-5)$ thì $\cos(\vec{a},\vec{b})=-\dfrac{8\sqrt{42}}{63}$.}{s}
  \y{Với $\vec{a}=(0;0;3)$, $\vec{b}=(-1;-4;-5)$ thì $(-\vec{a}+3\vec{b})\cdot(\vec{a}+3\vec{b})=369$.}{d}
  \y{$\vec{a}\cdot\vec{b}=-16$.}{s}}

\cauDS[id=C2-B8-D08-TH-050,muc=TH,loigiai={Kiểm tra từng mệnh đề bằng công thức tọa độ. a) Đúng. b) Sai. c) Đúng. d) Sai.}]{Trong không gian $Oxyz$, cho $\vec{a}=(-3;0;-1)$, $\vec{b}=(0;-3;5)$. Xét tính đúng sai của các mệnh đề sau:}{%
  \y{$\vec{a}\cdot\vec{b}=-5$.}{d}
  \y{Với $\vec{a}=(-3;0;-1)$, $\vec{b}=(0;-3;5)$ thì $(\vec{a}+2\vec{b})\cdot(3\vec{a}+3\vec{b})=567$.}{s}
  \y{Với $A(-1;0;-5)$, $B(0;-4;-4)$, $C(1;0;4)$ thì $\cos\widehat{ABC}=\dfrac{7\sqrt{2}}{54}$.}{d}
  \y{Với $\vec{a}=(-3;0;-1)$, $\vec{b}=(0;-3;5)$ thì $(3\vec{a}+2\vec{b})\cdot(\vec{a}+\vec{b})=70$.}{s}}

\cauDS[id=C2-B8-D08-TH-051,muc=TH,loigiai={Kiểm tra từng mệnh đề bằng công thức tọa độ. a) Đúng. b) Đúng. c) Đúng. d) Sai.}]{Trong không gian $Oxyz$, cho $\vec{a}=(0;5;-2)$, $\vec{b}=(4;4;2)$. Xét tính đúng sai của các mệnh đề sau:}{%
  \y{Với $\vec{a}=(0;5;-2)$, $\vec{b}=(4;4;2)$ thì $(2\vec{a}+\vec{b})\cdot(2\vec{a}+\vec{b})=216$.}{d}
  \y{Với $\vec{a}=(0;5;-2)$, $\vec{b}=(4;4;2)$ thì $\cos(\vec{a},\vec{b})=\dfrac{8\sqrt{29}}{87}$.}{d}
  \y{$\vec{a}\cdot\vec{b}=16$.}{d}
  \y{Với $\vec{a}=(0;5;-2)$, $\vec{b}=(4;4;2)$ thì $(\vec{a}+2\vec{b})\cdot(3\vec{a}+\vec{b})=543$.}{s}}

\cauDS[id=C2-B8-D08-TH-052,muc=TH,loigiai={Kiểm tra từng mệnh đề bằng công thức tọa độ. a) Đúng. b) Đúng. c) Sai. d) Đúng.}]{Trong không gian $Oxyz$, cho $\vec{a}=(0;5;-3)$, $\vec{b}=(2;4;-1)$. Xét tính đúng sai của các mệnh đề sau:}{%
  \y{Với $\vec{a}=(0;5;-3)$, $\vec{b}=(2;4;-1)$ thì $(-\vec{a}+2\vec{b})\cdot(\vec{a}+\vec{b})=31$.}{d}
  \y{Với $A(-3;2;-1)$, $B(3;-1;-2)$, $C(4;3;-3)$ thì $\cos\widehat{ABC}=\dfrac{5\sqrt{23}}{138}$.}{d}
  \y{$\vec{a}\cdot\vec{b}=20$.}{s}
  \y{Với $\vec{a}=(0;5;-3)$, $\vec{b}=(2;4;-1)$ thì $\cos(\vec{a},\vec{b})=\dfrac{23\sqrt{714}}{714}$.}{d}}

\cauDS[id=C2-B8-D08-TH-053,muc=TH,loigiai={Kiểm tra từng mệnh đề bằng công thức tọa độ. a) Đúng. b) Đúng. c) Đúng. d) Sai.}]{Trong không gian $Oxyz$, cho $\vec{a}=(-2;-3;-1)$, $\vec{b}=(-3;-4;4)$. Xét tính đúng sai của các mệnh đề sau:}{%
  \y{Với $\vec{a}=(-2;-3;-1)$, $\vec{b}=(-3;-4;4)$ thì $(-\vec{a}+3\vec{b})\cdot(\vec{a}+3\vec{b})=355$.}{d}
  \y{$\vec{a}\cdot\vec{b}=14$.}{d}
  \y{Với $A(-1;-2;5)$, $B(5;-2;-3)$, $C(4;-5;-2)$ thì $\cos\widehat{ABC}=\dfrac{7\sqrt{11}}{55}$.}{d}
  \y{Với $\vec{a}=(-2;-3;-1)$, $\vec{b}=(-3;-4;4)$ thì $(2\vec{a}+2\vec{b})\cdot(\vec{a}+2\vec{b})=275$.}{s}}

\cauDS[id=C2-B8-D08-TH-054,muc=TH,loigiai={Kiểm tra từng mệnh đề bằng công thức tọa độ. a) Đúng. b) Sai. c) Đúng. d) Đúng.}]{Trong không gian $Oxyz$, cho $\vec{a}=(-4;2;5)$, $\vec{b}=(4;2;2)$. Xét tính đúng sai của các mệnh đề sau:}{%
  \y{$\vec{a}\cdot\vec{b}=-2$.}{d}
  \y{Với $\vec{a}=(-4;2;5)$, $\vec{b}=(4;2;2)$ thì $\cos(\vec{a},\vec{b})=-\dfrac{\sqrt{30}}{180}$.}{s}
  \y{Với $\vec{a}=(-4;2;5)$, $\vec{b}=(4;2;2)$ thì $(3\vec{a}+2\vec{b})\cdot(2\vec{a}-\vec{b})=220$.}{d}
  \y{Với $\vec{a}=(-4;2;5)$, $\vec{b}=(4;2;2)$ thì $(2\vec{a}+\vec{b})\cdot(3\vec{a}+2\vec{b})=304$.}{d}}

\cauDS[id=C2-B8-D08-TH-055,muc=TH,loigiai={Kiểm tra từng mệnh đề bằng công thức tọa độ. a) Đúng. b) Đúng. c) Sai. d) Sai.}]{Trong không gian $Oxyz$, cho $\vec{a}=(-1;-1;5)$, $\vec{b}=(-1;-5;5)$. Xét tính đúng sai của các mệnh đề sau:}{%
  \y{Với $A(-1;3;-2)$, $B(4;3;0)$, $C(0;-5;-4)$ thì $\cos\widehat{ABC}=\dfrac{7\sqrt{174}}{174}$.}{d}
  \y{$\vec{a}\cdot\vec{b}=31$.}{d}
  \y{Với $A(-1;3;-2)$, $B(4;3;0)$, $C(0;-5;-4)$ thì $\cos\widehat{ABC}=\dfrac{9\sqrt{174}}{232}$.}{s}
  \y{Với $\vec{a}=(-1;-1;5)$, $\vec{b}=(-1;-5;5)$ thì $(2\vec{a}+2\vec{b})\cdot(3\vec{a}+\vec{b})=511$.}{s}}

\cauDS[id=C2-B8-D08-TH-056,muc=TH,loigiai={Kiểm tra từng mệnh đề bằng công thức tọa độ. a) Đúng. b) Sai. c) Sai. d) Sai.}]{Trong không gian $Oxyz$, cho $\vec{a}=(0;1;-5)$, $\vec{b}=(-3;4;-2)$. Xét tính đúng sai của các mệnh đề sau:}{%
  \y{Với $A(-1;-2;1)$, $B(4;3;0)$, $C(2;-3;-1)$ thì $\cos\widehat{ABC}=\dfrac{13\sqrt{2091}}{697}$.}{d}
  \y{Với $\vec{a}=(0;1;-5)$, $\vec{b}=(-3;4;-2)$ thì $(\vec{a}+2\vec{b})\cdot(-\vec{a}+3\vec{b})=325$.}{s}
  \y{Với $\vec{a}=(0;1;-5)$, $\vec{b}=(-3;4;-2)$ thì $(\vec{a}+3\vec{b})\cdot(2\vec{a}+\vec{b})=238$.}{s}
  \y{$\vec{a}\cdot\vec{b}=16$.}{s}}

\cauDS[id=C2-B8-D08-TH-057,muc=TH,loigiai={Kiểm tra từng mệnh đề bằng công thức tọa độ. a) Đúng. b) Sai. c) Sai. d) Đúng.}]{Trong không gian $Oxyz$, cho $\vec{a}=(-5;-5;-4)$, $\vec{b}=(-3;1;-2)$. Xét tính đúng sai của các mệnh đề sau:}{%
  \y{$\vec{a}\cdot\vec{b}=18$.}{d}
  \y{Với $\vec{a}=(-5;-5;-4)$, $\vec{b}=(-3;1;-2)$ thì $\cos(\vec{a},\vec{b})=\dfrac{17\sqrt{231}}{462}$.}{s}
  \y{$\vec{a}\cdot\vec{b}=37$.}{s}
  \y{Với $\vec{a}=(-5;-5;-4)$, $\vec{b}=(-3;1;-2)$ thì $(3\vec{a}+2\vec{b})\cdot(-\vec{a}+2\vec{b})=-70$.}{d}}

\cauDS[id=C2-B8-D08-TH-058,muc=TH,loigiai={Kiểm tra từng mệnh đề bằng công thức tọa độ. a) Sai. b) Sai. c) Sai. d) Đúng.}]{Xét tính đúng sai của các mệnh đề sau:}{%
  \y{Với $\vec{a}=(-5;3;-1)$, $\vec{b}=(-4;-3;0)$ thì $(3\vec{a}+3\vec{b})\cdot(-\vec{a}+\vec{b})=-59$.}{s}
  \y{Với $A(1;1;4)$, $B(-3;4;1)$, $C(0;-2;-5)$ thì $\cos\widehat{ABC}=-\dfrac{2\sqrt{34}}{51}$.}{s}
  \y{Với $A(1;1;4)$, $B(-3;4;1)$, $C(0;-2;-5)$ thì $\cos\widehat{ABC}=\dfrac{11\sqrt{34}}{306}$.}{s}
  \y{Với $\vec{a}=(-5;3;-1)$, $\vec{b}=(-4;-3;0)$ thì $\cos(\vec{a},\vec{b})=\dfrac{11\sqrt{35}}{175}$.}{d}}

\cauDS[id=C2-B8-D08-TH-059,muc=TH,loigiai={Kiểm tra từng mệnh đề bằng công thức tọa độ. a) Đúng. b) Sai. c) Sai. d) Đúng.}]{Xét tính đúng sai của các mệnh đề sau:}{%
  \y{Với $\vec{a}=(1;1;-4)$, $\vec{b}=(-2;-5;-1)$ thì $(2\vec{a}+2\vec{b})\cdot(2\vec{a}+\vec{b})=114$.}{d}
  \y{Với $A(4;-3;-5)$, $B(4;1;-1)$, $C(-4;-2;5)$ thì $\cos\widehat{ABC}=-\dfrac{13\sqrt{218}}{872}$.}{s}
  \y{Với $A(4;-3;-5)$, $B(4;1;-1)$, $C(-4;-2;5)$ thì $\cos\widehat{ABC}=\dfrac{3\sqrt{218}}{218}$.}{s}
  \y{Với $\vec{a}=(1;1;-4)$, $\vec{b}=(-2;-5;-1)$ thì $\cos(\vec{a},\vec{b})=-\dfrac{\sqrt{15}}{30}$.}{d}}

\cauDS[id=C2-B8-D08-TH-060,muc=TH,loigiai={Kiểm tra từng mệnh đề bằng công thức tọa độ. a) Sai. b) Đúng. c) Đúng. d) Đúng.}]{Xét tính đúng sai của các mệnh đề sau:}{%
  \y{Với $A(5;5;-3)$, $B(-3;-5;1)$, $C(1;-4;-4)$ thì $\cos\widehat{ABC}=\dfrac{\sqrt{210}}{20}$.}{s}
  \y{Với $A(5;5;-3)$, $B(-3;-5;1)$, $C(1;-4;-4)$ thì $\cos\widehat{ABC}=\dfrac{31\sqrt{210}}{630}$.}{d}
  \y{Với $\vec{a}=(-4;5;5)$, $\vec{b}=(2;2;2)$ thì $\cos(\vec{a},\vec{b})=\dfrac{\sqrt{22}}{11}$.}{d}
  \y{Với $\vec{a}=(-4;5;5)$, $\vec{b}=(2;2;2)$ thì $(3\vec{a}-\vec{b})\cdot(-\vec{a}+3\vec{b})=-114$.}{d}}

\cauDS[id=C2-B8-D08-TH-061,muc=TH,loigiai={Kiểm tra từng mệnh đề bằng công thức tọa độ. a) Sai. b) Đúng. c) Đúng. d) Sai.}]{Trong không gian $Oxyz$, cho $\vec{a}=(2;1;-1)$, $\vec{b}=(1;3;3)$. Xét tính đúng sai của các mệnh đề sau:}{%
  \y{Với $\vec{a}=(2;1;-1)$, $\vec{b}=(1;3;3)$ thì $\cos(\vec{a},\vec{b})=\dfrac{\sqrt{114}}{114}$.}{s}
  \y{Với $\vec{a}=(2;1;-1)$, $\vec{b}=(1;3;3)$ thì $(2\vec{a}+2\vec{b})\cdot(2\vec{a}+3\vec{b})=158$.}{d}
  \y{$\vec{a}\cdot\vec{b}=2$.}{d}
  \y{Với $\vec{a}=(2;1;-1)$, $\vec{b}=(1;3;3)$ thì $\cos(\vec{a},\vec{b})=-\dfrac{\sqrt{114}}{57}$.}{s}}

\cauDS[id=C2-B8-D08-TH-062,muc=TH,loigiai={Kiểm tra từng mệnh đề bằng công thức tọa độ. a) Đúng. b) Sai. c) Đúng. d) Đúng.}]{Trong không gian $Oxyz$, cho $\vec{a}=(-4;1;1)$, $\vec{b}=(4;-3;5)$. Xét tính đúng sai của các mệnh đề sau:}{%
  \y{$\vec{a}\cdot\vec{b}=-14$.}{d}
  \y{Với $A(-1;-3;0)$, $B(-2;2;0)$, $C(-3;-5;3)$ thì $\cos\widehat{ABC}=\dfrac{35\sqrt{1534}}{1534}$.}{s}
  \y{Với $\vec{a}=(-4;1;1)$, $\vec{b}=(4;-3;5)$ thì $\cos(\vec{a},\vec{b})=-\dfrac{7}{15}$.}{d}
  \y{Với $\vec{a}=(-4;1;1)$, $\vec{b}=(4;-3;5)$ thì $(-\vec{a}-\vec{b})\cdot(-\vec{a}-\vec{b})=40$.}{d}}

\cauDS[id=C2-B8-D08-TH-063,muc=TH,loigiai={Kiểm tra từng mệnh đề bằng công thức tọa độ. a) Đúng. b) Đúng. c) Sai. d) Đúng.}]{Trong không gian $Oxyz$, cho $\vec{a}=(-5;2;-1)$, $\vec{b}=(-3;-5;1)$. Xét tính đúng sai của các mệnh đề sau:}{%
  \y{$\vec{a}\cdot\vec{b}=4$.}{d}
  \y{Với $\vec{a}=(-5;2;-1)$, $\vec{b}=(-3;-5;1)$ thì $\cos(\vec{a},\vec{b})=\dfrac{2\sqrt{42}}{105}$.}{d}
  \y{Với $\vec{a}=(-5;2;-1)$, $\vec{b}=(-3;-5;1)$ thì $(-\vec{a}+2\vec{b})\cdot(\vec{a}+3\vec{b})=179$.}{s}
  \y{Với $A(0;-5;1)$, $B(4;-4;-5)$, $C(1;5;2)$ thì $\cos\widehat{ABC}=\dfrac{45\sqrt{7367}}{7367}$.}{d}}

\cauDS[id=C2-B8-D08-TH-064,muc=TH,loigiai={Kiểm tra từng mệnh đề bằng công thức tọa độ. a) Sai. b) Đúng. c) Sai. d) Sai.}]{Xét tính đúng sai của các mệnh đề sau:}{%
  \y{Với $A(-1;-2;-4)$, $B(-2;0;3)$, $C(0;2;-1)$ thì $\cos\widehat{ABC}=\dfrac{25}{36}$.}{s}
  \y{Với $\vec{a}=(3;2;5)$, $\vec{b}=(-4;-4;-1)$ thì $\cos(\vec{a},\vec{b})=-\dfrac{25\sqrt{1254}}{1254}$.}{d}
  \y{Với $\vec{a}=(3;2;5)$, $\vec{b}=(-4;-4;-1)$ thì $(2\vec{a}+3\vec{b})\cdot(-\vec{a}+\vec{b})=24$.}{s}
  \y{Với $A(-1;-2;-4)$, $B(-2;0;3)$, $C(0;2;-1)$ thì $\cos\widehat{ABC}=-\dfrac{13}{18}$.}{s}}

\cauTLN[id=C2-B8-D08-TH-065,muc=TH,loigiai={Ta có $-\vec{a}+\vec{b}=(-2;1;1)$, $3\vec{a}+3\vec{b}=(-24;-15;21)$ nên tích vô hướng bằng $54$.}]{Cho $\vec{a}=(-3;-3;3)$, $\vec{b}=(-5;-2;4)$. Tính $(-\vec{a}+\vec{b})\cdot(3\vec{a}+3\vec{b})$.}{54}

\cauTLN[id=C2-B8-D08-TH-066,muc=TH,loigiai={Ta có $-\vec{a}+\vec{b}=(2;0;-6)$, $2\vec{a}+2\vec{b}=(-4;4;8)$ nên tích vô hướng bằng $-56$.}]{Cho $\vec{a}=(-2;1;5)$, $\vec{b}=(0;1;-1)$. Tính $(-\vec{a}+\vec{b})\cdot(2\vec{a}+2\vec{b})$.}{-56}

\cauTLN[id=C2-B8-D08-TH-067,muc=TH,loigiai={Ta có $\vec{a}+3\vec{b}=(7;15;9)$, $3\vec{a}+2\vec{b}=(14;10;-1)$ nên tích vô hướng bằng $239$.}]{Cho $\vec{a}=(4;0;-3)$, $\vec{b}=(1;5;4)$. Tính $(\vec{a}+3\vec{b})\cdot(3\vec{a}+2\vec{b})$.}{239}

\cauTLN[id=C2-B8-D08-TH-068,muc=TH,loigiai={Ta có $\vec{a}-\vec{b}=(6;1;3)$, $2\vec{a}+\vec{b}=(0;2;0)$ nên tích vô hướng bằng $2$.}]{Cho $\vec{a}=(2;1;1)$, $\vec{b}=(-4;0;-2)$. Tính $(\vec{a}-\vec{b})\cdot(2\vec{a}+\vec{b})$.}{2}

\cauTLN[id=C2-B8-D08-TH-069,muc=TH,loigiai={Ta có $3\vec{a}-\vec{b}=(-4;-20;-17)$, $3\vec{a}+\vec{b}=(-14;-10;-13)$ nên tích vô hướng bằng $477$.}]{Cho $\vec{a}=(-3;-5;-5)$, $\vec{b}=(-5;5;2)$. Tính $(3\vec{a}-\vec{b})\cdot(3\vec{a}+\vec{b})$.}{477}

\cauTLN[id=C2-B8-D08-TH-070,muc=TH,loigiai={Ta có $-\vec{a}+3\vec{b}=(3;12;-9)$, $-\vec{a}+3\vec{b}=(3;12;-9)$ nên tích vô hướng bằng $234$.}]{Cho $\vec{a}=(0;-3;0)$, $\vec{b}=(1;3;-3)$. Tính $(-\vec{a}+3\vec{b})\cdot(-\vec{a}+3\vec{b})$.}{234}

\cauTLN[id=C2-B8-D08-TH-071,muc=TH,loigiai={Ta có $\vec{a}+3\vec{b}=(-18;-17;13)$, $\vec{a}+2\vec{b}=(-13;-12;10)$ nên tích vô hướng bằng $568$.}]{Cho $\vec{a}=(-3;-2;4)$, $\vec{b}=(-5;-5;3)$. Tính $(\vec{a}+3\vec{b})\cdot(\vec{a}+2\vec{b})$.}{568}

\cauTLN[id=C2-B8-D08-TH-072,muc=TH,loigiai={Ta có $\vec{a}+\vec{b}=(-6;0;1)$, $2\vec{a}-\vec{b}=(-3;3;-1)$ nên tích vô hướng bằng $17$.}]{Cho $\vec{a}=(-3;1;0)$, $\vec{b}=(-3;-1;1)$. Tính $(\vec{a}+\vec{b})\cdot(2\vec{a}-\vec{b})$.}{17}

\cauTLN[id=C2-B8-D08-TH-073,muc=TH,loigiai={Ta có $\vec{a}\cdot\vec{b}=3\cdot0+5\cdot2-1\cdot-4=14$.}]{Trong không gian $Oxyz$, cho $\vec{a}=(3;5;-1)$, $\vec{b}=(0;2;-4)$. Tính $\vec{a}\cdot\vec{b}$.}{14}

\cauTLN[id=C2-B8-D08-TH-074,muc=TH,loigiai={Ta có $\vec{a}\cdot\vec{b}=1\cdot4-2\cdot-5-4\cdot4=-2$.}]{Trong không gian $Oxyz$, cho $\vec{a}=(1;-2;-4)$, $\vec{b}=(4;-5;4)$. Tính $\vec{a}\cdot\vec{b}$.}{-2}

\cauTLN[id=C2-B8-D08-TH-075,muc=TH,loigiai={Ta có $\vec{a}\cdot\vec{b}=-5\cdot-1+0\cdot3-4\cdot-3=17$.}]{Trong không gian $Oxyz$, cho $\vec{a}=(-5;0;-4)$, $\vec{b}=(-1;3;-3)$. Tính $\vec{a}\cdot\vec{b}$.}{17}

\cauTLN[id=C2-B8-D08-TH-076,muc=TH,loigiai={Ta có $\vec{a}\cdot\vec{b}=-4\cdot3-3\cdot1+4\cdot5=5$.}]{Trong không gian $Oxyz$, cho $\vec{a}=(-4;-3;4)$, $\vec{b}=(3;1;5)$. Tính $\vec{a}\cdot\vec{b}$.}{5}

\cauTLN[id=C2-B8-D08-TH-077,muc=TH,loigiai={Ta có $\vec{a}\cdot\vec{b}=0\cdot-2+2\cdot3-4\cdot4=-10$.}]{Trong không gian $Oxyz$, cho $\vec{a}=(0;2;-4)$, $\vec{b}=(-2;3;4)$. Tính $\vec{a}\cdot\vec{b}$.}{-10}

\cauTLN[id=C2-B8-D08-TH-078,muc=TH,loigiai={Ta có $\vec{a}\cdot\vec{b}=-5\cdot3+1\cdot0-4\cdot3=-27$.}]{Trong không gian $Oxyz$, cho $\vec{a}=(-5;1;-4)$, $\vec{b}=(3;0;3)$. Tính $\vec{a}\cdot\vec{b}$.}{-27}

\cauTLN[id=C2-B8-D08-TH-079,muc=TH,loigiai={Ta có $\vec{a}\cdot\vec{b}=4\cdot5-1\cdot2+1\cdot-5=13$.}]{Trong không gian $Oxyz$, cho $\vec{a}=(4;-1;1)$, $\vec{b}=(5;2;-5)$. Tính $\vec{a}\cdot\vec{b}$.}{13}

\cauTLN[id=C2-B8-D08-TH-080,muc=TH,loigiai={Ta có $\vec{a}\cdot\vec{b}=0\cdot4-1\cdot2+5\cdot3=13$.}]{Trong không gian $Oxyz$, cho $\vec{a}=(0;-1;5)$, $\vec{b}=(4;2;3)$. Tính $\vec{a}\cdot\vec{b}$.}{13}

\cauTN[id=C2-B8-D08-VD-001,muc=VD,loigiai={Gọi $M(0;0;t)$. $\overrightarrow{MA}\cdot\overrightarrow{MB}=0$ cho $t^2-(-3)t+(2)=0$, có hai nghiệm $t=-2$ và $t=-1$.}]{Điểm $M$ thuộc trục $Oz$ sao cho $\overrightarrow{MA}\perp\overrightarrow{MB}$, với $A(3;-1;-1)$, $B(0;0;-2)$. Tọa độ của $M$ là}
  {$M(0;0;-2)\text{ hoặc }M(0;0;-1)$.}
  {$M(0;0;-2)\text{ hoặc }M(0;0;0)$.}{$M(0;0;1)\text{ hoặc }M(0;0;2)$.}{$M(0;0;-2)\text{ hoặc }M(0;0;1)$.}

\cauTN[id=C2-B8-D08-VD-002,muc=VD,loigiai={Gọi $M(0;0;t)$. $\overrightarrow{MA}\cdot\overrightarrow{MB}=0$ cho $t^2-(5)t+(6)=0$, có hai nghiệm $t=2$ và $t=3$.}]{Điểm $M$ thuộc trục $Oz$ sao cho $\overrightarrow{MA}\perp\overrightarrow{MB}$, với $A(1;-3;2)$, $B(0;0;3)$. Tọa độ của $M$ là}
  {$M(0;0;2)\text{ hoặc }M(0;0;3)$.}
  {$M(0;0;2)\text{ hoặc }M(0;0;4)$.}{$M(0;0;-3)\text{ hoặc }M(0;0;-2)$.}{$M(0;0;-3)\text{ hoặc }M(0;0;2)$.}

\cauTN[id=C2-B8-D08-VD-003,muc=VD,loigiai={Gọi $M(0;0;t)$. $\overrightarrow{MA}\cdot\overrightarrow{MB}=0$ cho $t^2-(1)t+(-6)=0$, có hai nghiệm $t=-2$ và $t=3$.}]{Điểm $M$ thuộc trục $Oz$ sao cho $\overrightarrow{MA}\perp\overrightarrow{MB}$, với $A(2;-3;1)$, $B(-3;0;0)$. Tọa độ của $M$ là}
  {$M(0;0;-2)\text{ hoặc }M(0;0;3)$.}
  {$M(0;0;-2)\text{ hoặc }M(0;0;4)$.}{$M(0;0;-1)\text{ hoặc }M(0;0;3)$.}{$M(0;0;-3)\text{ hoặc }M(0;0;2)$.}

\cauTN[id=C2-B8-D08-VD-004,muc=VD,loigiai={Gọi $M(0;0;t)$. $\overrightarrow{MA}\cdot\overrightarrow{MB}=0$ cho $t^2-(0)t+(-4)=0$, có hai nghiệm $t=-2$ và $t=2$.}]{Điểm $M$ thuộc trục $Oz$ sao cho $\overrightarrow{MA}\perp\overrightarrow{MB}$, với $A(0;0;2)$, $B(2;0;-2)$. Tọa độ của $M$ là}
  {$M(0;0;-2)\text{ hoặc }M(0;0;2)$.}
  {$M(0;0;-2)\text{ hoặc }M(0;0;3)$.}{$M(0;0;-1)\text{ hoặc }M(0;0;2)$.}{$M(0;0;-1)\text{ hoặc }M(0;0;3)$.}

\cauTN[id=C2-B8-D08-VD-005,muc=VD,loigiai={Gọi $M(0;0;t)$. $\overrightarrow{MA}\cdot\overrightarrow{MB}=0$ cho $t^2-(2)t+(0)=0$, có hai nghiệm $t=0$ và $t=2$.}]{Điểm $M$ thuộc trục $Oz$ sao cho $\overrightarrow{MA}\perp\overrightarrow{MB}$, với $A(0;-1;1)$, $B(2;1;1)$. Tọa độ của $M$ là}
  {$M(0;0;0)\text{ hoặc }M(0;0;2)$.}
  {$M(0;0;0)\text{ hoặc }M(0;0;3)$.}{$M(0;0;1)\text{ hoặc }M(0;0;2)$.}{$M(0;0;-2)\text{ hoặc }M(0;0;0)$.}

\cauTN[id=C2-B8-D08-VD-006,muc=VD,loigiai={Gọi $M(0;0;t)$. $\overrightarrow{MA}\cdot\overrightarrow{MB}=0$ cho $t^2-(1)t+(-12)=0$, có hai nghiệm $t=-3$ và $t=4$.}]{Điểm $M$ thuộc trục $Oz$ sao cho $\overrightarrow{MA}\perp\overrightarrow{MB}$, với $A(3;0;-2)$, $B(-2;1;3)$. Tọa độ của $M$ là}
  {$M(0;0;-3)\text{ hoặc }M(0;0;4)$.}
  {$M(0;0;-3)\text{ hoặc }M(0;0;5)$.}{$M(0;0;-2)\text{ hoặc }M(0;0;4)$.}{$M(0;0;-4)\text{ hoặc }M(0;0;3)$.}

\cauTN[id=C2-B8-D08-VD-007,muc=VD,loigiai={Gọi $M(0;0;t)$. $\overrightarrow{MA}\cdot\overrightarrow{MB}=0$ cho $t^2-(-3)t+(2)=0$, có hai nghiệm $t=-2$ và $t=-1$.}]{Điểm $M$ thuộc trục $Oz$ sao cho $\overrightarrow{MA}\perp\overrightarrow{MB}$, với $A(-1;1;-3)$, $B(1;3;0)$. Tọa độ của $M$ là}
  {$M(0;0;-2)\text{ hoặc }M(0;0;-1)$.}
  {$M(0;0;-2)\text{ hoặc }M(0;0;0)$.}{$M(0;0;1)\text{ hoặc }M(0;0;2)$.}{$M(0;0;-2)\text{ hoặc }M(0;0;1)$.}

\cauTN[id=C2-B8-D08-VD-008,muc=VD,loigiai={Gọi $M(0;0;t)$. $\overrightarrow{MA}\cdot\overrightarrow{MB}=0$ cho $t^2-(2)t+(0)=0$, có hai nghiệm $t=0$ và $t=2$.}]{Điểm $M$ thuộc trục $Oz$ sao cho $\overrightarrow{MA}\perp\overrightarrow{MB}$, với $A(2;3;2)$, $B(3;-2;0)$. Tọa độ của $M$ là}
  {$M(0;0;0)\text{ hoặc }M(0;0;2)$.}
  {$M(0;0;0)\text{ hoặc }M(0;0;3)$.}{$M(0;0;1)\text{ hoặc }M(0;0;2)$.}{$M(0;0;-2)\text{ hoặc }M(0;0;0)$.}

\cauTN[id=C2-B8-D08-VD-009,muc=VD,loigiai={Gọi $M(0;0;t)$. $\overrightarrow{MA}\cdot\overrightarrow{MB}=0$ cho $t^2-(-1)t+(-6)=0$, có hai nghiệm $t=-3$ và $t=2$.}]{Điểm $M$ thuộc trục $Oz$ sao cho $\overrightarrow{MA}\perp\overrightarrow{MB}$, với $A(2;2;-3)$, $B(-1;1;2)$. Tọa độ của $M$ là}
  {$M(0;0;-3)\text{ hoặc }M(0;0;2)$.}
  {$M(0;0;-3)\text{ hoặc }M(0;0;3)$.}{$M(0;0;-2)\text{ hoặc }M(0;0;2)$.}{$M(0;0;-2)\text{ hoặc }M(0;0;3)$.}

\cauTN[id=C2-B8-D08-VD-010,muc=VD,loigiai={Gọi $M(0;0;t)$. $\overrightarrow{MA}\cdot\overrightarrow{MB}=0$ cho $t^2-(2)t+(0)=0$, có hai nghiệm $t=0$ và $t=2$.}]{Điểm $M$ thuộc trục $Oz$ sao cho $\overrightarrow{MA}\perp\overrightarrow{MB}$, với $A(-1;-1;1)$, $B(-1;2;1)$. Tọa độ của $M$ là}
  {$M(0;0;0)\text{ hoặc }M(0;0;2)$.}
  {$M(0;0;0)\text{ hoặc }M(0;0;3)$.}{$M(0;0;1)\text{ hoặc }M(0;0;2)$.}{$M(0;0;-2)\text{ hoặc }M(0;0;0)$.}

\cauTN[id=C2-B8-D08-VD-011,muc=VD,loigiai={Gọi $M(0;0;t)$. $\overrightarrow{MA}\cdot\overrightarrow{MB}=0$ cho $t^2-(4)t+(3)=0$, có hai nghiệm $t=1$ và $t=3$.}]{Điểm $M$ thuộc trục $Oz$ sao cho $\overrightarrow{MA}\perp\overrightarrow{MB}$, với $A(3;-2;2)$, $B(1;2;2)$. Tọa độ của $M$ là}
  {$M(0;0;1)\text{ hoặc }M(0;0;3)$.}
  {$M(0;0;1)\text{ hoặc }M(0;0;4)$.}{$M(0;0;2)\text{ hoặc }M(0;0;3)$.}{$M(0;0;-3)\text{ hoặc }M(0;0;-1)$.}

\cauTN[id=C2-B8-D08-VD-012,muc=VD,loigiai={Gọi $M(0;0;t)$. $\overrightarrow{MA}\cdot\overrightarrow{MB}=0$ cho $t^2-(1)t+(0)=0$, có hai nghiệm $t=0$ và $t=1$.}]{Điểm $M$ thuộc trục $Oz$ sao cho $\overrightarrow{MA}\perp\overrightarrow{MB}$, với $A(-2;2;3)$, $B(-3;0;-2)$. Tọa độ của $M$ là}
  {$M(0;0;0)\text{ hoặc }M(0;0;1)$.}
  {$M(0;0;0)\text{ hoặc }M(0;0;2)$.}{$M(0;0;-1)\text{ hoặc }M(0;0;0)$.}{$M(0;0;1)\text{ hoặc }M(0;0;2)$.}

\cauTN[id=C2-B8-D08-VD-013,muc=VD,loigiai={Ta có $\overrightarrow{BA}=(3;-6;-10)$, $\overrightarrow{BC}=(-4;-10;-6)$ nên $\cos\widehat{ABC}=\dfrac{\overrightarrow{BA}\cdot\overrightarrow{BC}}{|\overrightarrow{BA}||\overrightarrow{BC}|}=\dfrac{27\sqrt{5510}}{2755}$.}]{Cho tam giác $ABC$ với $A(1;-2;-6)$, $B(-2;4;4)$, $C(-6;-6;-2)$. Giá trị $\cos\widehat{ABC}$ bằng}
  {$\dfrac{27\sqrt{5510}}{2755}$.}
  {$\dfrac{109\sqrt{5510}}{11020}$.}{$-\dfrac{27\sqrt{5510}}{2755}$.}{$\dfrac{107\sqrt{5510}}{11020}$.}

\cauTN[id=C2-B8-D08-VD-014,muc=VD,loigiai={Ta có $\overrightarrow{BA}=(3;0;1)$, $\overrightarrow{BC}=(-1;11;1)$ nên $\cos\widehat{ABC}=\dfrac{\overrightarrow{BA}\cdot\overrightarrow{BC}}{|\overrightarrow{BA}||\overrightarrow{BC}|}=-\dfrac{\sqrt{1230}}{615}$.}]{Cho tam giác $ABC$ với $A(-1;-5;-2)$, $B(-4;-5;-3)$, $C(-5;6;-2)$. Giá trị $\cos\widehat{ABC}$ bằng}
  {$-\dfrac{\sqrt{1230}}{615}$.}
  {$-\dfrac{\sqrt{1230}}{1230}$.}{$\dfrac{\sqrt{1230}}{615}$.}{$-\dfrac{\sqrt{1230}}{410}$.}

\cauTN[id=C2-B8-D08-VD-015,muc=VD,loigiai={Ta có $\overrightarrow{BA}=(0;10;0)$, $\overrightarrow{BC}=(7;5;-1)$ nên $\cos\widehat{ABC}=\dfrac{\overrightarrow{BA}\cdot\overrightarrow{BC}}{|\overrightarrow{BA}||\overrightarrow{BC}|}=\dfrac{\sqrt{3}}{3}$.}]{Cho tam giác $ABC$ với $A(-6;5;2)$, $B(-6;-5;2)$, $C(1;0;1)$. Giá trị $\cos\widehat{ABC}$ bằng}
  {$\dfrac{\sqrt{3}}{3}$.}
  {$\dfrac{17\sqrt{3}}{50}$.}{$-\dfrac{\sqrt{3}}{3}$.}{$\dfrac{49\sqrt{3}}{150}$.}

\cauTN[id=C2-B8-D08-VD-016,muc=VD,loigiai={Ta có $\overrightarrow{BA}=(4;6;-6)$, $\overrightarrow{BC}=(-3;-1;-8)$ nên $\cos\widehat{ABC}=\dfrac{\overrightarrow{BA}\cdot\overrightarrow{BC}}{|\overrightarrow{BA}||\overrightarrow{BC}|}=\dfrac{15\sqrt{407}}{814}$.}]{Cho tam giác $ABC$ với $A(2;5;0)$, $B(-2;-1;6)$, $C(-5;-2;-2)$. Giá trị $\cos\widehat{ABC}$ bằng}
  {$\dfrac{15\sqrt{407}}{814}$.}
  {$\dfrac{31\sqrt{407}}{1628}$.}{$-\dfrac{15\sqrt{407}}{814}$.}{$\dfrac{29\sqrt{407}}{1628}$.}

\cauTN[id=C2-B8-D08-VD-017,muc=VD,loigiai={Ta có $\overrightarrow{BA}=(0;5;-7)$, $\overrightarrow{BC}=(3;7;0)$ nên $\cos\widehat{ABC}=\dfrac{\overrightarrow{BA}\cdot\overrightarrow{BC}}{|\overrightarrow{BA}||\overrightarrow{BC}|}=\dfrac{35\sqrt{1073}}{2146}$.}]{Cho tam giác $ABC$ với $A(0;4;-1)$, $B(0;-1;6)$, $C(3;6;6)$. Giá trị $\cos\widehat{ABC}$ bằng}
  {$\dfrac{35\sqrt{1073}}{2146}$.}
  {$\dfrac{18\sqrt{1073}}{1073}$.}{$-\dfrac{35\sqrt{1073}}{2146}$.}{$\dfrac{17\sqrt{1073}}{1073}$.}

\cauTN[id=C2-B8-D08-VD-018,muc=VD,loigiai={Ta có $\overrightarrow{BA}=(6;0;-10)$, $\overrightarrow{BC}=(2;0;0)$ nên $\cos\widehat{ABC}=\dfrac{\overrightarrow{BA}\cdot\overrightarrow{BC}}{|\overrightarrow{BA}||\overrightarrow{BC}|}=\dfrac{3\sqrt{34}}{34}$.}]{Cho tam giác $ABC$ với $A(3;-2;-5)$, $B(-3;-2;5)$, $C(-1;-2;5)$. Giá trị $\cos\widehat{ABC}$ bằng}
  {$\dfrac{3\sqrt{34}}{34}$.}
  {$\dfrac{13\sqrt{34}}{136}$.}{$-\dfrac{3\sqrt{34}}{34}$.}{$\dfrac{11\sqrt{34}}{136}$.}

\cauTN[id=C2-B8-D08-VD-019,muc=VD,loigiai={Ta có $\overrightarrow{BA}=(-2;2;-1)$, $\overrightarrow{BC}=(3;-2;-1)$ nên $\cos\widehat{ABC}=\dfrac{\overrightarrow{BA}\cdot\overrightarrow{BC}}{|\overrightarrow{BA}||\overrightarrow{BC}|}=-\dfrac{3\sqrt{14}}{14}$.}]{Cho tam giác $ABC$ với $A(-4;0;2)$, $B(-2;-2;3)$, $C(1;-4;2)$. Giá trị $\cos\widehat{ABC}$ bằng}
  {$-\dfrac{3\sqrt{14}}{14}$.}
  {$-\dfrac{4\sqrt{14}}{21}$.}{$\dfrac{3\sqrt{14}}{14}$.}{$-\dfrac{5\sqrt{14}}{21}$.}

\cauTN[id=C2-B8-D08-VD-020,muc=VD,loigiai={Ta có $\overrightarrow{BA}=(4;4;-7)$, $\overrightarrow{BC}=(10;0;-1)$ nên $\cos\widehat{ABC}=\dfrac{\overrightarrow{BA}\cdot\overrightarrow{BC}}{|\overrightarrow{BA}||\overrightarrow{BC}|}=\dfrac{47\sqrt{101}}{909}$.}]{Cho tam giác $ABC$ với $A(-2;0;-4)$, $B(-6;-4;3)$, $C(4;-4;2)$. Giá trị $\cos\widehat{ABC}$ bằng}
  {$\dfrac{47\sqrt{101}}{909}$.}
  {$\dfrac{16\sqrt{101}}{303}$.}{$-\dfrac{47\sqrt{101}}{909}$.}{$\dfrac{46\sqrt{101}}{909}$.}

\cauTN[id=C2-B8-D08-VD-021,muc=VD,loigiai={Ta có $\overrightarrow{BA}=(2;-1;-3)$, $\overrightarrow{BC}=(-2;2;3)$ nên $\cos\widehat{ABC}=\dfrac{\overrightarrow{BA}\cdot\overrightarrow{BC}}{|\overrightarrow{BA}||\overrightarrow{BC}|}=-\dfrac{15\sqrt{238}}{238}$.}]{Cho tam giác $ABC$ với $A(5;-3;0)$, $B(3;-2;3)$, $C(1;0;6)$. Giá trị $\cos\widehat{ABC}$ bằng}
  {$-\dfrac{15\sqrt{238}}{238}$.}
  {$-\dfrac{\sqrt{238}}{17}$.}{$\dfrac{15\sqrt{238}}{238}$.}{$-\dfrac{13\sqrt{238}}{238}$.}

\cauTN[id=C2-B8-D08-VD-022,muc=VD,loigiai={Ta có $\overrightarrow{BA}=(-3;3;-1)$, $\overrightarrow{BC}=(3;-3;1)$ nên $\cos\widehat{ABC}=\dfrac{\overrightarrow{BA}\cdot\overrightarrow{BC}}{|\overrightarrow{BA}||\overrightarrow{BC}|}=-1$.}]{Cho tam giác $ABC$ với $A(-6;5;4)$, $B(-3;2;5)$, $C(0;-1;6)$. Giá trị $\cos\widehat{ABC}$ bằng}
  {$-1$.}
  {$-\dfrac{18}{19}$.}{$1$.}{$-\dfrac{17}{19}$.}

\cauTN[id=C2-B8-D08-VD-023,muc=VD,loigiai={Ta có $\overrightarrow{BA}=(6;-1;-3)$, $\overrightarrow{BC}=(9;-8;-1)$ nên $\cos\widehat{ABC}=\dfrac{\overrightarrow{BA}\cdot\overrightarrow{BC}}{|\overrightarrow{BA}||\overrightarrow{BC}|}=\dfrac{65\sqrt{1679}}{3358}$.}]{Cho tam giác $ABC$ với $A(0;1;-6)$, $B(-6;2;-3)$, $C(3;-6;-4)$. Giá trị $\cos\widehat{ABC}$ bằng}
  {$\dfrac{65\sqrt{1679}}{3358}$.}
  {$\dfrac{33\sqrt{1679}}{1679}$.}{$-\dfrac{65\sqrt{1679}}{3358}$.}{$\dfrac{32\sqrt{1679}}{1679}$.}

\cauTN[id=C2-B8-D08-VD-024,muc=VD,loigiai={Ta có $\overrightarrow{BA}=(5;-9;2)$, $\overrightarrow{BC}=(-4;-9;-1)$ nên $\cos\widehat{ABC}=\dfrac{\overrightarrow{BA}\cdot\overrightarrow{BC}}{|\overrightarrow{BA}||\overrightarrow{BC}|}=\dfrac{59\sqrt{55}}{770}$.}]{Cho tam giác $ABC$ với $A(5;-6;5)$, $B(0;3;3)$, $C(-4;-6;2)$. Giá trị $\cos\widehat{ABC}$ bằng}
  {$\dfrac{59\sqrt{55}}{770}$.}
  {$\dfrac{6\sqrt{55}}{77}$.}{$-\dfrac{59\sqrt{55}}{770}$.}{$\dfrac{29\sqrt{55}}{385}$.}

\cauTN[id=C2-B8-D08-VD-025,muc=VD,loigiai={Ta có $-\vec{a}+\vec{b}=(0;11;-9)$, $2\vec{a}+\vec{b}=(-6;-4;3)$ nên tích vô hướng bằng $-71$.}]{Cho $\vec{a}=(-2;-5;4)$, $\vec{b}=(-2;6;-5)$. Tích vô hướng $(-\vec{a}+\vec{b})\cdot(2\vec{a}+\vec{b})$ bằng}
  {$-71$.}
  {$-70$.}{$-142$.}{$-213$.}

\cauTN[id=C2-B8-D08-VD-026,muc=VD,loigiai={Ta có $2\vec{a}+3\vec{b}=(1;12;-10)$, $3\vec{a}+2\vec{b}=(-1;3;-15)$ nên tích vô hướng bằng $185$.}]{Cho $\vec{a}=(-1;-3;-5)$, $\vec{b}=(1;6;0)$. Tích vô hướng $(2\vec{a}+3\vec{b})\cdot(3\vec{a}+2\vec{b})$ bằng}
  {$185$.}
  {$188$.}{$371$.}{$183$.}

\cauTN[id=C2-B8-D08-VD-027,muc=VD,loigiai={Ta có $3\vec{a}-\vec{b}=(7;13;-18)$, $-\vec{a}-\vec{b}=(-5;-7;2)$ nên tích vô hướng bằng $-162$.}]{Cho $\vec{a}=(3;5;-5)$, $\vec{b}=(2;2;3)$. Tích vô hướng $(3\vec{a}-\vec{b})\cdot(-\vec{a}-\vec{b})$ bằng}
  {$-162$.}
  {$-164$.}{$-486$.}{$-324$.}

\cauTN[id=C2-B8-D08-VD-028,muc=VD,loigiai={Ta có $2\vec{a}+\vec{b}=(-11;4;-6)$, $\vec{a}+\vec{b}=(-8;0;0)$ nên tích vô hướng bằng $88$.}]{Cho $\vec{a}=(-3;4;-6)$, $\vec{b}=(-5;-4;6)$. Tích vô hướng $(2\vec{a}+\vec{b})\cdot(\vec{a}+\vec{b})$ bằng}
  {$88$.}
  {$44$.}{$-88$.}{$87$.}

\cauTN[id=C2-B8-D08-VD-029,muc=VD,loigiai={Ta có $2\vec{a}-\vec{b}=(-9;3;8)$, $2\vec{a}-\vec{b}=(-9;3;8)$ nên tích vô hướng bằng $154$.}]{Cho $\vec{a}=(-2;-1;5)$, $\vec{b}=(5;-5;2)$. Tích vô hướng $(2\vec{a}-\vec{b})\cdot(2\vec{a}-\vec{b})$ bằng}
  {$154$.}
  {$153$.}{$152$.}{$155$.}

\cauTN[id=C2-B8-D08-VD-030,muc=VD,loigiai={Ta có $-\vec{a}+2\vec{b}=(0;8;-11)$, $2\vec{a}+2\vec{b}=(-18;20;-2)$ nên tích vô hướng bằng $182$.}]{Cho $\vec{a}=(-6;4;3)$, $\vec{b}=(-3;6;-4)$. Tích vô hướng $(-\vec{a}+2\vec{b})\cdot(2\vec{a}+2\vec{b})$ bằng}
  {$182$.}
  {$180$.}{$185$.}{$181$.}

\cauTN[id=C2-B8-D08-VD-031,muc=VD,loigiai={Ta có $-\vec{a}+\vec{b}=(7;-2;-5)$, $2\vec{a}+3\vec{b}=(-4;-21;0)$ nên tích vô hướng bằng $14$.}]{Cho $\vec{a}=(-5;-3;3)$, $\vec{b}=(2;-5;-2)$. Tích vô hướng $(-\vec{a}+\vec{b})\cdot(2\vec{a}+3\vec{b})$ bằng}
  {$14$.}
  {$12$.}{$16$.}{$-14$.}

\cauTN[id=C2-B8-D08-VD-032,muc=VD,loigiai={Ta có $2\vec{a}+\vec{b}=(2;4;-6)$, $3\vec{a}-\vec{b}=(3;16;-19)$ nên tích vô hướng bằng $184$.}]{Cho $\vec{a}=(1;4;-5)$, $\vec{b}=(0;-4;4)$. Tích vô hướng $(2\vec{a}+\vec{b})\cdot(3\vec{a}-\vec{b})$ bằng}
  {$184$.}
  {$185$.}{$552$.}{$368$.}

\cauTN[id=C2-B8-D08-VD-033,muc=VD,loigiai={Ta có $2\vec{a}+\vec{b}=(-12;0;13)$, $2\vec{a}+3\vec{b}=(-24;-8;23)$ nên tích vô hướng bằng $587$.}]{Cho $\vec{a}=(-3;2;4)$, $\vec{b}=(-6;-4;5)$. Tích vô hướng $(2\vec{a}+\vec{b})\cdot(2\vec{a}+3\vec{b})$ bằng}
  {$587$.}
  {$586$.}{$\dfrac{587}{2}$.}{$588$.}

\cauTN[id=C2-B8-D08-VD-034,muc=VD,loigiai={Ta có $\vec{a}+2\vec{b}=(-6;-2;6)$, $\vec{a}+3\vec{b}=(-12;-3;8)$ nên tích vô hướng bằng $126$.}]{Cho $\vec{a}=(6;0;2)$, $\vec{b}=(-6;-1;2)$. Tích vô hướng $(\vec{a}+2\vec{b})\cdot(\vec{a}+3\vec{b})$ bằng}
  {$126$.}
  {$-126$.}{$124$.}{$125$.}

\cauTN[id=C2-B8-D08-VD-035,muc=VD,loigiai={Ta có $2\vec{a}-\vec{b}=(-6;-7;-1)$, $\vec{a}+\vec{b}=(-12;-5;7)$ nên tích vô hướng bằng $100$.}]{Cho $\vec{a}=(-6;-4;2)$, $\vec{b}=(-6;-1;5)$. Tích vô hướng $(2\vec{a}-\vec{b})\cdot(\vec{a}+\vec{b})$ bằng}
  {$100$.}
  {$201$.}{$200$.}{$101$.}

\cauTN[id=C2-B8-D08-VD-036,muc=VD,loigiai={Ta có $3\vec{a}+\vec{b}=(-12;-2;23)$, $2\vec{a}+\vec{b}=(-9;0;17)$ nên tích vô hướng bằng $499$.}]{Cho $\vec{a}=(-3;-2;6)$, $\vec{b}=(-3;4;5)$. Tích vô hướng $(3\vec{a}+\vec{b})\cdot(2\vec{a}+\vec{b})$ bằng}
  {$499$.}
  {$500$.}{$501$.}{$\dfrac{499}{2}$.}

\cauDS[id=C2-B8-D08-VD-037,muc=VD,loigiai={Kiểm tra từng mệnh đề bằng công thức tọa độ. a) Đúng. b) Đúng. c) Sai. d) Đúng.}]{Xét tính đúng sai của các mệnh đề sau:}{%
  \y{Với $A(-3;4;3)$, $B(-1;0;-2)$, $C(-3;-1;0)$ thì $\cos\widehat{ABC}=\dfrac{2\sqrt{5}}{9}$.}{d}
  \y{Với $\vec{a}=(1;-1;5)$, $\vec{b}=(3;3;-1)$ thì $(3\vec{a}+2\vec{b})\cdot(3\vec{a}+\vec{b})=236$.}{d}
  \y{Với $\vec{a}=(1;-1;5)$, $\vec{b}=(3;3;-1)$ thì $(-\vec{a}-\vec{b})\cdot(2\vec{a}+3\vec{b})=-171$.}{s}
  \y{Với $\vec{a}=(1;-1;5)$, $\vec{b}=(3;3;-1)$ thì $\cos(\vec{a},\vec{b})=-\dfrac{5\sqrt{57}}{171}$.}{d}}

\cauDS[id=C2-B8-D08-VD-038,muc=VD,loigiai={Kiểm tra từng mệnh đề bằng công thức tọa độ. a) Đúng. b) Sai. c) Đúng. d) Sai.}]{Xét tính đúng sai của các mệnh đề sau:}{%
  \y{Với $\vec{a}=(6;-5;3)$, $\vec{b}=(5;3;1)$ thì $(\vec{a}+3\vec{b})\cdot(3\vec{a}+2\vec{b})=618$.}{d}
  \y{Với $\vec{a}=(6;-5;3)$, $\vec{b}=(5;3;1)$ thì $\cos(\vec{a},\vec{b})=\dfrac{19\sqrt{2}}{70}$.}{s}
  \y{Với $\vec{a}=(6;-5;3)$, $\vec{b}=(5;3;1)$ thì $(-\vec{a}+3\vec{b})\cdot(2\vec{a}+\vec{b})=55$.}{d}
  \y{Với $A(2;6;4)$, $B(0;-6;0)$, $C(0;0;-3)$ thì $\cos\widehat{ABC}=\dfrac{59\sqrt{205}}{1230}$.}{s}}

\cauDS[id=C2-B8-D08-VD-039,muc=VD,loigiai={Kiểm tra từng mệnh đề bằng công thức tọa độ. a) Đúng. b) Đúng. c) Đúng. d) Sai.}]{Xét tính đúng sai của các mệnh đề sau:}{%
  \y{Với $A(5;2;5)$, $B(0;2;5)$, $C(2;-5;-3)$ thì $\cos\widehat{ABC}=\dfrac{2\sqrt{13}}{39}$.}{d}
  \y{Với $\vec{a}=(-3;-4;6)$, $\vec{b}=(-5;4;-5)$ thì $(\vec{a}+3\vec{b})\cdot(-\vec{a}+2\vec{b})=366$.}{d}
  \y{Với $\vec{a}=(-3;-4;6)$, $\vec{b}=(-5;4;-5)$ thì $\cos(\vec{a},\vec{b})=-\dfrac{31\sqrt{4026}}{4026}$.}{d}
  \y{Với $\vec{a}=(-3;-4;6)$, $\vec{b}=(-5;4;-5)$ thì $(3\vec{a}-\vec{b})\cdot(2\vec{a}+\vec{b})=538$.}{s}}

\cauDS[id=C2-B8-D08-VD-040,muc=VD,loigiai={Kiểm tra từng mệnh đề bằng công thức tọa độ. a) Đúng. b) Sai. c) Đúng. d) Sai.}]{Xét tính đúng sai của các mệnh đề sau:}{%
  \y{Với $\vec{a}=(6;-6;5)$, $\vec{b}=(-2;6;3)$ thì $\cos(\vec{a},\vec{b})=-\dfrac{33\sqrt{97}}{679}$.}{d}
  \y{Với $A(5;4;-6)$, $B(0;-6;-2)$, $C(2;5;4)$ thì $\cos\widehat{ABC}=\dfrac{97\sqrt{22701}}{22701}$.}{s}
  \y{Với $A(5;4;-6)$, $B(0;-6;-2)$, $C(2;5;4)$ thì $\cos\widehat{ABC}=\dfrac{32\sqrt{22701}}{7567}$.}{d}
  \y{Với $\vec{a}=(6;-6;5)$, $\vec{b}=(-2;6;3)$ thì $(2\vec{a}+3\vec{b})\cdot(3\vec{a}+3\vec{b})=1584$.}{s}}

\cauDS[id=C2-B8-D08-VD-041,muc=VD,loigiai={Kiểm tra từng mệnh đề bằng công thức tọa độ. a) Đúng. b) Đúng. c) Đúng. d) Sai.}]{Xét tính đúng sai của các mệnh đề sau:}{%
  \y{Với $A(-2;1;1)$, $B(6;3;4)$, $C(5;-1;6)$ thì $\cos\widehat{ABC}=\dfrac{10\sqrt{33}}{231}$.}{d}
  \y{Với $\vec{a}=(0;-3;6)$, $\vec{b}=(-2;-2;-2)$ thì $(3\vec{a}-\vec{b})\cdot(2\vec{a}+2\vec{b})=222$.}{d}
  \y{Với $\vec{a}=(0;-3;6)$, $\vec{b}=(-2;-2;-2)$ thì $\cos(\vec{a},\vec{b})=-\dfrac{\sqrt{15}}{15}$.}{d}
  \y{Với $\vec{a}=(0;-3;6)$, $\vec{b}=(-2;-2;-2)$ thì $(3\vec{a}+3\vec{b})\cdot(3\vec{a}+2\vec{b})=1161$.}{s}}

\cauDS[id=C2-B8-D08-VD-042,muc=VD,loigiai={Kiểm tra từng mệnh đề bằng công thức tọa độ. a) Đúng. b) Sai. c) Sai. d) Đúng.}]{Xét tính đúng sai của các mệnh đề sau:}{%
  \y{Với $\vec{a}=(-3;-5;0)$, $\vec{b}=(-4;-5;-4)$ thì $(\vec{a}+\vec{b})\cdot(-\vec{a}+\vec{b})=23$.}{d}
  \y{Với $A(-5;6;-2)$, $B(-2;-3;-6)$, $C(-2;3;2)$ thì $\cos\widehat{ABC}=\dfrac{87\sqrt{106}}{1060}$.}{s}
  \y{Với $\vec{a}=(-3;-5;0)$, $\vec{b}=(-4;-5;-4)$ thì $\cos(\vec{a},\vec{b})=-\dfrac{37\sqrt{1938}}{1938}$.}{s}
  \y{Với $A(-5;6;-2)$, $B(-2;-3;-6)$, $C(-2;3;2)$ thì $\cos\widehat{ABC}=\dfrac{43\sqrt{106}}{530}$.}{d}}

\cauDS[id=C2-B8-D08-VD-043,muc=VD,loigiai={Kiểm tra từng mệnh đề bằng công thức tọa độ. a) Đúng. b) Sai. c) Sai. d) Đúng.}]{Xét tính đúng sai của các mệnh đề sau:}{%
  \y{Với $\vec{a}=(-5;-3;0)$, $\vec{b}=(5;1;5)$ thì $(\vec{a}+2\vec{b})\cdot(2\vec{a}+\vec{b})=30$.}{d}
  \y{Với $A(3;-1;2)$, $B(-2;-3;6)$, $C(-5;2;1)$ thì $\cos\widehat{ABC}=-\dfrac{\sqrt{295}}{59}$.}{s}
  \y{Với $\vec{a}=(-5;-3;0)$, $\vec{b}=(5;1;5)$ thì $\cos(\vec{a},\vec{b})=-\dfrac{29\sqrt{6}}{102}$.}{s}
  \y{Với $\vec{a}=(-5;-3;0)$, $\vec{b}=(5;1;5)$ thì $(\vec{a}+\vec{b})\cdot(3\vec{a}+2\vec{b})=64$.}{d}}

\cauDS[id=C2-B8-D08-VD-044,muc=VD,loigiai={Kiểm tra từng mệnh đề bằng công thức tọa độ. a) Đúng. b) Sai. c) Đúng. d) Sai.}]{Xét tính đúng sai của các mệnh đề sau:}{%
  \y{Với $A(3;-4;0)$, $B(-1;4;-2)$, $C(-2;-3;2)$ thì $\cos\widehat{ABC}=\dfrac{5\sqrt{154}}{77}$.}{d}
  \y{Với $\vec{a}=(3;1;-6)$, $\vec{b}=(5;-5;4)$ thì $(2\vec{a}+\vec{b})\cdot(\vec{a}+3\vec{b})=193$.}{s}
  \y{Với $\vec{a}=(3;1;-6)$, $\vec{b}=(5;-5;4)$ thì $\cos(\vec{a},\vec{b})=-\dfrac{7\sqrt{759}}{759}$.}{d}
  \y{Với $\vec{a}=(3;1;-6)$, $\vec{b}=(5;-5;4)$ thì $\cos(\vec{a},\vec{b})=\dfrac{7\sqrt{759}}{759}$.}{s}}

\cauDS[id=C2-B8-D08-VD-045,muc=VD,loigiai={Kiểm tra từng mệnh đề bằng công thức tọa độ. a) Đúng. b) Sai. c) Đúng. d) Sai.}]{Xét tính đúng sai của các mệnh đề sau:}{%
  \y{Với $\vec{a}=(4;3;-2)$, $\vec{b}=(6;-4;5)$ thì $\cos(\vec{a},\vec{b})=\dfrac{2\sqrt{2233}}{2233}$.}{d}
  \y{Với $A(4;-6;5)$, $B(2;-6;-5)$, $C(3;-3;-4)$ thì $\cos\widehat{ABC}=\dfrac{\sqrt{286}}{52}$.}{s}
  \y{Với $\vec{a}=(4;3;-2)$, $\vec{b}=(6;-4;5)$ thì $(\vec{a}+3\vec{b})\cdot(2\vec{a}+2\vec{b})=536$.}{d}
  \y{Với $\vec{a}=(4;3;-2)$, $\vec{b}=(6;-4;5)$ thì $(2\vec{a}+\vec{b})\cdot(2\vec{a}+2\vec{b})=280$.}{s}}

\cauDS[id=C2-B8-D08-VD-046,muc=VD,loigiai={Kiểm tra từng mệnh đề bằng công thức tọa độ. a) Đúng. b) Đúng. c) Sai. d) Sai.}]{Xét tính đúng sai của các mệnh đề sau:}{%
  \y{Với $A(2;-5;3)$, $B(-4;6;-5)$, $C(5;-1;-3)$ thì $\cos\widehat{ABC}=\dfrac{147\sqrt{29614}}{29614}$.}{d}
  \y{Với $\vec{a}=(3;-4;1)$, $\vec{b}=(2;-6;4)$ thì $(3\vec{a}-\vec{b})\cdot(-\vec{a}+2\vec{b})=48$.}{d}
  \y{Với $\vec{a}=(3;-4;1)$, $\vec{b}=(2;-6;4)$ thì $(\vec{a}+\vec{b})\cdot(\vec{a}+\vec{b})=152$.}{s}
  \y{Với $\vec{a}=(3;-4;1)$, $\vec{b}=(2;-6;4)$ thì $\cos(\vec{a},\vec{b})=\dfrac{33\sqrt{91}}{364}$.}{s}}

\cauDS[id=C2-B8-D08-VD-047,muc=VD,loigiai={Kiểm tra từng mệnh đề bằng công thức tọa độ. a) Đúng. b) Đúng. c) Sai. d) Đúng.}]{Xét tính đúng sai của các mệnh đề sau:}{%
  \y{Với $\vec{a}=(0;-6;0)$, $\vec{b}=(-6;5;-3)$ thì $\cos(\vec{a},\vec{b})=-\dfrac{\sqrt{70}}{14}$.}{d}
  \y{Với $\vec{a}=(0;-6;0)$, $\vec{b}=(-6;5;-3)$ thì $(-\vec{a}-\vec{b})\cdot(2\vec{a}+\vec{b})=-52$.}{d}
  \y{Với $\vec{a}=(0;-6;0)$, $\vec{b}=(-6;5;-3)$ thì $(\vec{a}-\vec{b})\cdot(2\vec{a}+\vec{b})=29$.}{s}
  \y{Với $A(3;-6;6)$, $B(0;-5;2)$, $C(1;4;0)$ thì $\cos\widehat{ABC}=-\dfrac{7\sqrt{559}}{559}$.}{d}}

\cauDS[id=C2-B8-D08-VD-048,muc=VD,loigiai={Kiểm tra từng mệnh đề bằng công thức tọa độ. a) Đúng. b) Sai. c) Đúng. d) Sai.}]{Xét tính đúng sai của các mệnh đề sau:}{%
  \y{Với $A(2;-1;4)$, $B(-5;-1;5)$, $C(3;2;4)$ thì $\cos\widehat{ABC}=\dfrac{57\sqrt{37}}{370}$.}{d}
  \y{Với $A(2;-1;4)$, $B(-5;-1;5)$, $C(3;2;4)$ thì $\cos\widehat{ABC}=\dfrac{28\sqrt{37}}{185}$.}{s}
  \y{Với $\vec{a}=(6;0;6)$, $\vec{b}=(-6;-2;-6)$ thì $(-\vec{a}+3\vec{b})\cdot(-\vec{a}+3\vec{b})=1188$.}{d}
  \y{Với $\vec{a}=(6;0;6)$, $\vec{b}=(-6;-2;-6)$ thì $\cos(\vec{a},\vec{b})=-\dfrac{73\sqrt{38}}{456}$.}{s}}

\cauTLN[id=C2-B8-D08-VD-049,muc=VD,loigiai={Gọi $M(0;0;t)$. $\overrightarrow{MA}\cdot\overrightarrow{MB}=0$ cho $t^2-(0)t+(-4)=0$, có hai nghiệm $t=-2$ và $t=2$.}]{Điểm $M$ thuộc trục $Oz$ sao cho $\overrightarrow{MA}\perp\overrightarrow{MB}$, với $A(2;-2;2)$, $B(-3;-3;-2)$. Tính tổng các cao độ của tất cả các điểm $M$ thoả mãn.}{0}

\cauTLN[id=C2-B8-D08-VD-050,muc=VD,loigiai={Gọi $M(0;0;t)$. $\overrightarrow{MA}\cdot\overrightarrow{MB}=0$ cho $t^2-(2)t+(0)=0$, có hai nghiệm $t=0$ và $t=2$.}]{Điểm $M$ thuộc trục $Oz$ sao cho $\overrightarrow{MA}\perp\overrightarrow{MB}$, với $A(-1;-1;0)$, $B(-1;1;2)$. Tính tổng các cao độ của tất cả các điểm $M$ thoả mãn.}{2}

\cauTLN[id=C2-B8-D08-VD-051,muc=VD,loigiai={Gọi $M(0;0;t)$. $\overrightarrow{MA}\cdot\overrightarrow{MB}=0$ cho $t^2-(-1)t+(-12)=0$, có hai nghiệm $t=-4$ và $t=3$.}]{Điểm $M$ thuộc trục $Oz$ sao cho $\overrightarrow{MA}\perp\overrightarrow{MB}$, với $A(0;2;-3)$, $B(-3;-3;2)$. Tính tổng các cao độ của tất cả các điểm $M$ thoả mãn.}{-1}

\cauTLN[id=C2-B8-D08-VD-052,muc=VD,loigiai={Gọi $M(0;0;t)$. $\overrightarrow{MA}\cdot\overrightarrow{MB}=0$ cho $t^2-(2)t+(-3)=0$, có hai nghiệm $t=-1$ và $t=3$.}]{Điểm $M$ thuộc trục $Oz$ sao cho $\overrightarrow{MA}\perp\overrightarrow{MB}$, với $A(3;1;2)$, $B(-2;3;0)$. Tính tổng các cao độ của tất cả các điểm $M$ thoả mãn.}{2}

\cauTLN[id=C2-B8-D08-VD-053,muc=VD,loigiai={Gọi $M(0;0;t)$. $\overrightarrow{MA}\cdot\overrightarrow{MB}=0$ cho $t^2-(-3)t+(0)=0$, có hai nghiệm $t=-3$ và $t=0$.}]{Điểm $M$ thuộc trục $Oz$ sao cho $\overrightarrow{MA}\perp\overrightarrow{MB}$, với $A(2;3;-3)$, $B(-3;2;0)$. Tính tổng các cao độ của tất cả các điểm $M$ thoả mãn.}{-3}

\cauTLN[id=C2-B8-D08-VD-054,muc=VD,loigiai={Gọi $M(0;0;t)$. $\overrightarrow{MA}\cdot\overrightarrow{MB}=0$ cho $t^2-(-3)t+(0)=0$, có hai nghiệm $t=-3$ và $t=0$.}]{Điểm $M$ thuộc trục $Oz$ sao cho $\overrightarrow{MA}\perp\overrightarrow{MB}$, với $A(2;2;0)$, $B(-2;2;-3)$. Tính tổng các cao độ của tất cả các điểm $M$ thoả mãn.}{-3}

\cauTLN[id=C2-B8-D08-VD-055,muc=VD,loigiai={Ta có $2\vec{a}+3\vec{b}=(18;-2;-26)$, $\vec{a}+3\vec{b}=(18;2;-22)$ nên tích vô hướng bằng $892$.}]{Cho $\vec{a}=(0;-4;-4)$, $\vec{b}=(6;2;-6)$. Tính $(2\vec{a}+3\vec{b})\cdot(\vec{a}+3\vec{b})$.}{892}

\cauTLN[id=C2-B8-D08-VD-056,muc=VD,loigiai={Ta có $2\vec{a}+2\vec{b}=(8;-10;18)$, $2\vec{a}+\vec{b}=(2;-8;14)$ nên tích vô hướng bằng $348$.}]{Cho $\vec{a}=(-2;-3;5)$, $\vec{b}=(6;-2;4)$. Tính $(2\vec{a}+2\vec{b})\cdot(2\vec{a}+\vec{b})$.}{348}

\cauTLN[id=C2-B8-D08-VD-057,muc=VD,loigiai={Ta có $\vec{a}+2\vec{b}=(10;-5;6)$, $\vec{a}+\vec{b}=(5;-3;4)$ nên tích vô hướng bằng $89$.}]{Cho $\vec{a}=(0;-1;2)$, $\vec{b}=(5;-2;2)$. Tính $(\vec{a}+2\vec{b})\cdot(\vec{a}+\vec{b})$.}{89}

\cauTLN[id=C2-B8-D08-VD-058,muc=VD,loigiai={Ta có $2\vec{a}+2\vec{b}=(-6;-8;6)$, $3\vec{a}+3\vec{b}=(-9;-12;9)$ nên tích vô hướng bằng $204$.}]{Cho $\vec{a}=(1;-5;-2)$, $\vec{b}=(-4;1;5)$. Tính $(2\vec{a}+2\vec{b})\cdot(3\vec{a}+3\vec{b})$.}{204}

\cauTLN[id=C2-B8-D08-VD-059,muc=VD,loigiai={Ta có $-\vec{a}+3\vec{b}=(-4;-15;11)$, $-\vec{a}-\vec{b}=(4;-3;-5)$ nên tích vô hướng bằng $-26$.}]{Cho $\vec{a}=(-2;6;1)$, $\vec{b}=(-2;-3;4)$. Tính $(-\vec{a}+3\vec{b})\cdot(-\vec{a}-\vec{b})$.}{-26}

\cauTLN[id=C2-B8-D08-VD-060,muc=VD,loigiai={Ta có $2\vec{a}+3\vec{b}=(6;19;11)$, $-\vec{a}+\vec{b}=(-3;-2;2)$ nên tích vô hướng bằng $-34$.}]{Cho $\vec{a}=(3;5;1)$, $\vec{b}=(0;3;3)$. Tính $(2\vec{a}+3\vec{b})\cdot(-\vec{a}+\vec{b})$.}{-34}
